\documentclass[conference]{IEEEtran}

\usepackage{cite}
\usepackage{amsmath,amssymb}
\usepackage{booktabs}
\usepackage{graphicx}
\usepackage{enumitem}
\usepackage{url}
\usepackage{hyperref}

% IEEE standard description style: unboxed to prevent label-body overlap
\setlist[description]{
  style=unboxed,
  leftmargin=0pt,
  font=\normalfont\bfseries,
  itemsep=0.8ex plus 0.2ex minus 0.1ex,
  parsep=0pt
}

\title{Alert Aggregation from Early Warning Logs to Mitigate Alert Fatigue in Microservices:\\A Tool-Calling Agent Approach}

\author{
\IEEEauthorblockN{Dang Quang Hien}
\IEEEauthorblockA{
Faculty of Information Technology\\
University of Science, VNU-HCM\\
Ho Chi Minh City, Vietnam\\
dqhien@fit.hcmus.edu.vn
}
}

\begin{document}

\maketitle

\begin{abstract}
In distributed microservice architectures, cascading performance degradation during pre-failure phases routinely triggers alert storms, inundating on-call engineers with early-warning-level logs. Over 70\% of these alerts are noisy or unactionable, causing severe alert fatigue and delaying critical incident remediation. Existing alert management approaches primarily target post-mortem failure stages or rely on static rule-based grouping (e.g., Prometheus Alertmanager) and semantic similarity clustering, both of which lack awareness of dynamic service call topologies and causal propagation. This paper proposes a two-stage alert aggregation framework tailored specifically for early-warning telemetry in microservices. The framework integrates a lightweight deterministic fast filter for duplicate elimination and time-window pre-grouping with an LLM agent equipped with tool-calling capabilities (\texttt{get\_related\_alerts}, \texttt{get\_service\_dependency}, and \texttt{get\_runbook}). We evaluate on the RCAEval benchmark, which provides 735 failure cases across three systems (Online Boutique, Sock Shop, and Train Ticket), of which 359 cases expose raw log telemetry. Because RCAEval does not ship a severity column, we introduce an explicit three-tier severity assignment procedure that derives early-warning alerts from raw log messages and, for systems that emit no severity token, from faulty-versus-normal template frequency shifts. The proposed framework is designed to achieve an Alert Reduction Ratio (ARR) $\ge 60\%$ while strictly preserving the root cause service alert with a Root Cause Preservation Rate (RCPR) $\ge 95\%$. The primary contributions are: (1) formalizing a four-tier taxonomy of noisy early-warning alerts (duplicate, transient, unactionable, cascading), (2) introducing an agentic tool-calling mechanism for causal pre-failure alert triage, and (3) establishing a standardized multi-metric evaluation protocol balancing alert reduction against diagnostic safety.
\end{abstract}

\begin{IEEEkeywords}
Alert aggregation, microservices observability, early warning logs, alert fatigue, LLM agents, tool-calling, root cause preservation
\end{IEEEkeywords}

\section{Introduction}
Modern cloud computing has broadly adopted microservice architectures, decomposing monolithic systems into tens or hundreds of loosely coupled services interacting via network calls \cite{kuang2024cola, pham2025rcaeval}. While this architecture provides horizontal scalability and rapid deployment cycles, it fundamentally exacerbates operational complexity. Recent industry surveys reveal that enterprise operations manage an average of 101 distinct observability technologies, with 39\% of practitioners identifying cognitive overhead and informational complexity as their principal operational impediment \cite{grafana2025}. During anomalous episodes, cascading degradations—such as database connection exhaustion or downstream latency spikes—induce widespread symptom propagation across dependent services \cite{kuang2024cola, yu2025alertguardian}.

This propagation triggers an \textit{alert storm}, confronting Site Reliability Engineers (SREs) with hundreds of simultaneous notifications. Consequently, engineers suffer from \textit{alert fatigue}, a cognitive state in which desensitization leads to dismissed or overlooked genuine anomalies \cite{pagerduty2020, tariq2025alertfatigue, jalalvand2024alertprioritisation}. Unmanaged downtime carries staggering economic impacts: industry surveys indicate that over two-thirds of enterprise organizations incur outage costs exceeding \$300,000 per hour \cite{pagerduty2026, itic2024downtime}.

Crucially, current industrial alerting practices and academic literature suffer from two primary limitations:
\begin{enumerate}
    \item \textbf{Post-failure focus}: Existing Root Cause Analysis (RCA) techniques focus on post-mortem diagnosis after catastrophic failures (\texttt{ERROR} or \texttt{CRITICAL} state) have already violated Service Level Objectives (SLOs) \cite{sre-workbook-ch5, pham2025rcaeval}. Very few systems focus on the pre-failure degradation phase where telemetry manifests as \texttt{WARNING}-level logs \cite{mendes2021loglevels}.
    \item \textbf{Causal blindness in aggregation}: Traditional grouping mechanisms, such as Prometheus Alertmanager \cite{prometheus-alertmanager}, depend on static label matching and fixed time windows. Unsupervised semantic clustering (e.g., Sentence-BERT with HDBSCAN \cite{sbert, hdbscan}) groups syntactically similar messages but cannot discern topological dependencies, frequently collapsing unrelated service symptoms into misleading clusters.
\end{enumerate}

To bridge this gap, this paper introduces an agentic tool-calling framework for warning-level alert aggregation. The system utilizes deterministic fast filtering to prune immediate duplicates, followed by an LLM-driven agent that queries runtime topologies and synthetic runbooks to make causally coherent aggregation decisions.

\section{Related Work}
\subsection{Rule-Based and Label-Centric Alert Grouping}
Industrial alert pipelines rely predominantly on static rule engines, exemplified by Prometheus Alertmanager \cite{prometheus-alertmanager}. These systems apply label-matching configurations (\texttt{group\_by}) and temporal dampening intervals (\texttt{group\_wait}, \texttt{group\_interval}). While deterministic and computationally lightweight, rule-based systems are topology-blind; they cannot dynamically correlate alerts across upstream dependencies without manually written, fragile routing trees \cite{kuang2024cola}.

\subsection{Semantic and Machine Learning-Based Aggregation}
To overcome static rule fragility, data-driven approaches utilize log parsing techniques like Drain3 \cite{drain3} to extract log templates, subsequently generating vector representations via TF-IDF or Sentence-BERT \cite{sbert}. Unsupervised clustering algorithms such as HDBSCAN \cite{hdbscan} then cluster similar templates. While effective at lexical deduplication, semantic similarity alone is causality-blind: two independent services emitting timeout warnings will be conflated into a single group despite having completely disjoint root causes.

\subsection{Knowledge-Aware and LLM-Driven AIOps}
Recent advances explore topological graph correlation and hybrid LLM reasoning. COLA \cite{kuang2024cola} incorporates multi-source cloud knowledge for alert aggregation in large-scale cloud systems, while AlertGuardian \cite{yu2025alertguardian} implements life-cycle management across alert lifecycles. Furthermore, chain-of-thought prompting \cite{wei2022cot} has demonstrated strong diagnostic capabilities. Building upon these foundations, our approach addresses early warning logs specifically, employing constrained tool-calling to balance computational latency against causal reasoning accuracy.

\section{Methodology}
\subsection{Problem Formulation and Noise Taxonomy}
We define an early-warning alert stream as an ordered sequence $\mathcal{A} = \{a_1, a_2, \dots, a_N\}$, where each alert $a_i = \langle t_i, s_i, m_i, \ell_i \rangle$ consists of timestamp $t_i$, emitting service $s_i$, message template $m_i$, and severity tier $\ell_i$. In accordance with SRE alerting principles \cite{google-sre-book, google-sre-ch10, google-sre-ch11}, alerts must require immediate human intervention.

Because RCAEval does not expose a severity column, $\ell_i$ is \emph{assigned} by us rather than read from the dataset. We use a three-tier procedure:
\begin{enumerate}
    \item \textbf{Tier 1 (explicit).} For systems whose log framework emits a level token (Sock Shop and Train Ticket, both Spring Boot/log4j), parse $\ell_i$ from the message text itself; messages carrying \texttt{WARN} become \texttt{WARNING}, messages carrying \texttt{ERROR} or \texttt{SEVERE} become \texttt{ERROR}, and all others become \texttt{INFO}.
    \item \textbf{Tier 2 (frequency-inferred).} For systems that emit no level token at all (Online Boutique), a template $m$ is treated as anomalous when its frequency in the fault window exceeds its frequency in the normal window by a ratio $R$ (default $R \ge 3$) and its absolute frequency exceeds $N_{\min}$.
    \item \textbf{Tier 3 (ground-truth marker).} For the eight cases where RCAEval ships a \texttt{root\_cause.txt} indicator log line, that line is tagged \texttt{ROOT\_CAUSE} and used to compute an alert-level variant of RCPR.
\end{enumerate}

We further define two temporal windows from the injection timestamp $T_{\text{inject}}$, namely the normal window $[T_{\text{start}}, T_{\text{inject}})$ and the fault window $[T_{\text{inject}}, T_{\text{end}}]$. Note that the \texttt{normal\_timesteps} and \texttt{faulty\_timesteps} fields are integer counts of sample points, not timestamp lists, and therefore cannot be used directly for window assignment.

We formalize noisy alerts into four categories:
\begin{enumerate}
    \item \textbf{Duplicate Alert}: Multiple alerts where $s_i = s_j$, $m_i = m_j$, and $|t_i - t_j| \le \Delta t_{\text{dup}}$ (with $\Delta t_{\text{dup}} = 60\text{s}$).
    \item \textbf{Transient Alert}: An alert occurring during normal operating conditions ($t_i < T_{\text{inject}}$) whose template frequency ratio between fault and normal windows stays below a threshold $R_{\text{transient}}$, i.e., it subsides without administrative intervention.
    \item \textbf{Unactionable Alert}: A syntactically valid warning lacking an actionable remediation procedure or runbook mapping.
    \item \textbf{Cascading Alert}: An alert emitted by a service $s_i \neq s_{\text{root}}$ during the fault window ($t_i \ge T_{\text{inject}}$) that is reachable from $s_{\text{root}}$ in the service call graph, representing a propagated symptom rather than the originating fault.
\end{enumerate}

\subsection{Proposed Architecture: Tool-Calling Agent Framework}
The proposed pipeline comprises two consecutive stages:

\textbf{Stage 1: Deterministic Fast Filter.} Raw warning alerts first enter an online sliding-window deduplicator. Alerts sharing identical service identifiers and template hashes within a 60-second window are collapsed into a single candidate event, pruning an estimated 70--80\% of raw volume without invoking LLM inference.

\textbf{Stage 2: Tool-Augmented Reasoning Agent.} Unresolved candidate clusters are presented to an LLM agent configured with Chain-of-Thought prompting \cite{wei2022cot}. The agent is empowered with three deterministic environment tools:
\begin{itemize}
    \item \texttt{get\_related\_alerts(service, time\_window)}: Fetches concurrent alerts across neighboring services within a temporal radius $\Delta t$.
    \item \texttt{get\_service\_dependency(service)}: Queries the static microservice call graph $G=(V, E)$ to retrieve upstream dependencies and downstream dependents.
    \item \texttt{get\_runbook(fault\_type)}: Retrieves synthetic troubleshooting procedures and symptom-cause mappings for the candidate anomaly.
\end{itemize}
To guarantee operational bounded latency and avoid cyclic execution, the agent is constrained to a hard ceiling of at most two tool invocations per grouping decision.

\section{Experimental Setup}
\subsection{Benchmark Dataset}
Experiments utilize the public benchmark RCAEval \cite{pham2025rcaeval}, containing 735 annotated failure cases across three realistic microservice applications: Online Boutique (11 services), Sock Shop (13 services), and Train Ticket (40+ services). RCAEval provides multi-modal telemetry---logs, metrics, and distributed traces---accompanied by precise fault injection timestamps $T_{\text{inject}}$ and authoritative ground-truth root cause labels $s_{\text{root}}$.

A direct inspection of the benchmark revealed three properties that constrain the experimental design. First, only \textbf{359 of the 735 cases expose raw log telemetry}; the 375 cases of suite RE1 are metric-only. Second, each log file contains exactly three columns (\texttt{timestamp}, \texttt{container\_name}, \texttt{message}) and \textbf{no severity column}, so severity must be assigned by our three-tier procedure. Third, scanning all 49{,}652{,}771 log lines yields only \textbf{9{,}738 lines carrying a \texttt{WARN} token}, distributed highly unevenly: 9{,}662 lines originate from Sock Shop, 74 from Train Ticket, and only 2 from Online Boutique. Consequently, Sock Shop serves as the primary evaluation system for warning-tier aggregation, with Train Ticket retained for the cascading case study and Online Boutique restricted to the frequency-inferred tier. Cross-system validation is designated on the LEMMA-RCA dataset \cite{zheng2024lemma}, although we note that LEMMA-RCA aggregates log lines into three-dimensional time series whose golden-signal keywords cover only \texttt{error}, \texttt{exception}, and \texttt{critical}---not \texttt{warning}---so it does not provide an independent warning tier.

\subsection{Experimental Arms}
We benchmark the proposed framework across four comparative experimental arms:
\begin{itemize}
    \item \textbf{Arm 1 (Rule-Based Baseline)}: Prometheus Alertmanager simulation with standard label matching on service and severity, bounded by \texttt{group\_wait} (30s) and \texttt{group\_interval} (5m) \cite{prometheus-alertmanager}.
    \item \textbf{Arm 2 (Semantic Similarity)}: Drain3 log template extraction \cite{drain3} paired with Sentence-BERT embeddings (\texttt{all-MiniLM-L6-v2}) \cite{sbert} and HDBSCAN density clustering \cite{hdbscan}.
    \item \textbf{Arm 3 (Temporal-Spatial Correlation)}: Sliding window correlation ($\Delta t \in \{30, 60, 120\}\text{s}$) traversing the service call dependency graph with hop distance $k \le 2$. Because RCAEval does not ship the call graph as an artifact, we derive it from the distributed traces included for 240 of the 359 log-bearing cases, falling back to the published deployment manifests of the three open-source reference applications where trace coverage is absent.
    \item \textbf{Arm 4 (Tool-Calling Agent -- Proposed)}: Two-stage fast filter combined with the tool-calling LLM reasoning agent.
\end{itemize}

\subsection{Evaluation Metrics}
Evaluation relies on four primary quantitative metrics:
\begin{enumerate}
    \item \textbf{Alert Reduction Ratio (ARR)}: Measures overall volume compression:
    \begin{equation}
    \text{ARR} = \left(1 - \frac{N_{\text{groups}}}{N_{\text{raw\_alerts}}}\right) \times 100\%
    \end{equation}
    Target: $\text{ARR} \ge 60\%$.
    
    \item \textbf{Root Cause Preservation Rate (RCPR)}: Ensures safety by verifying that the primary incident cluster retains the genuine root cause evidence. We define RCPR at two granularities. The \emph{service-level} variant is the primary metric and spans all cases that expose logs ($M = 359$):
    \begin{equation}
    \text{RCPR}_{\text{svc}} = \frac{1}{M} \sum_{k=1}^{M} \mathbb{I}\left(s_{\text{root}}^{(k)} \in g_{\text{primary}}^{(k)}\right) \times 100\%
    \end{equation}
    The \emph{alert-level} variant is a secondary metric restricted to the $M'$ cases for which RCAEval ships a \texttt{root\_cause.txt} indicator line that is itself a warning-tier log ($M' = 4$):
    \begin{equation}
    \text{RCPR}_{\text{alert}} = \frac{1}{M'} \sum_{k=1}^{M'} \mathbb{I}\left(a_{\text{root}}^{(k)} \in g_{\text{primary}}^{(k)}\right) \times 100\%
    \end{equation}
    Target: $\text{RCPR} \ge 95\%$ (safety threshold).
    
    \item \textbf{Pairwise Grouping F1-Score ($F1_p$)}: Harmonic mean of pairwise precision ($P_p$) and pairwise recall ($R_p$) over alert pairs sharing identical root causes. Target: $F1_p \ge 0.85$.
    
    \item \textbf{Group Purity}: Average proportion of majority-cause alerts per aggregated group. Target: $\text{Purity} \ge 0.80$.
\end{enumerate}

\section{Expected Outcomes and Analysis Plan}
This section states the anticipated behaviour of each arm and the analysis plan. The figures below are \emph{research targets and expected trade-offs}, not measured results; no benchmark has been executed at the time of writing.

\begin{itemize}
    \item \textbf{Arm 1 (Rule-based)} is expected to achieve moderate compression ($\text{ARR}$ in the order of $40\%$) but comparatively low pairwise F1, because strict service-boundary isolation prevents downstream cascading symptoms from being aggregated.
    \item \textbf{Arm 2 (Semantic)} is expected to deliver aggressive reduction ($\text{ARR} > 70\%$) while degrading root cause preservation, erroneously collapsing lexically similar timeout logs across unrelated microservice boundaries. A concrete case in our data illustrates both failure modes: within Train Ticket, 74 warning lines concerning lost MongoDB connections normalise into 16 distinct templates that differ only by instance name, so lexical grouping yields 16 spurious clusters whereas semantic grouping should yield one.
    \item \textbf{Arm 3 (Temporal-Spatial)} is expected to capture cascading alerts along adjacent dependency edges but to remain brittle against fluctuating propagation delays. Because warning-tier cascading is scarce in RCAEval (74 lines, Train Ticket only), this arm will additionally be evaluated on an extended stream that includes the error tier, and the extension will be declared explicitly.
    \item \textbf{Arm 4 (Proposed Agent)} is expected to satisfy both optimisation objectives by corroborating dependency topology and runbook logic prior to group assignment.
\end{itemize}

Ablation analysis will quantify the marginal contribution of each tool, and the paired comparison across arms will be tested for significance (paired $t$-test with Wilcoxon signed-rank as a non-parametric fallback).

\section{Discussion}
\subsection{Operational Feasibility and Latency}
A primary challenge of LLM-based operations is inference latency and financial token expenditure. By filtering 70--80\% of high-frequency repetitive alerts at Stage 1, Stage 2 processes only structurally distinct candidates, constraining end-to-end processing latency to sub-second windows per incident window.

\subsection{Threats to Validity}
Internal validity is constrained by two factors specific to this benchmark. First, and most importantly, RCAEval carries \textbf{no severity column}, so the warning tier is not supplied by the dataset but \emph{derived} by our three-tier assignment procedure. Any conclusion is therefore conditional on that procedure; we mitigate the risk by reporting the explicit parsing rules and by validating the extracted templates against the eight cases for which RCAEval ships a ground-truth root cause indicator line. Second, synthetic fault injection patterns in RCAEval may exhibit cleaner separation than heterogeneous enterprise production logs.

External validity is constrained by the scarcity of warning-tier telemetry. Warning lines are concentrated in a single system (Sock Shop accounts for 99.2\% of them), and the only observed instance of warning-tier cascading consists of 74 lines in Train Ticket. Claims about cross-system generalisation of warning-tier aggregation therefore cannot be supported by this benchmark alone. We mitigate this by (a) evaluating cascading on an explicitly declared extended stream that includes the error tier, (b) reporting the frequency-inferred tier separately for Online Boutique, and (c) cross-evaluating against LEMMA-RCA \cite{zheng2024lemma}, while noting that LEMMA-RCA likewise defines no warning golden signal.

Construct validity is constrained by the fact that the \emph{unactionable} category depends on a synthetic runbook that we author ourselves, since RCAEval provides no SOP or runbook. We therefore report results both with and without the unactionable category.

\section{Conclusion}
This study presented a tool-calling LLM agent framework designed for early warning alert aggregation in microservices. By systematically targeting the pre-failure degradation phase and formalizing a four-category noise model, the framework addresses the limitations of topology-blind semantic models and rigid rule engines. Future work will investigate real-time streaming implementations and reinforcement learning feedback for dynamic tool policy optimization.

\appendix
\section*{Glossary of Terms}
\addcontentsline{toc}{section}{Glossary of Terms}
% =============================================================================
% GLOSSARY OF TERMS — Alert Aggregation in Microservices
% Standalone file for % =============================================================================
% GLOSSARY OF TERMS — Alert Aggregation in Microservices
% Standalone file for % =============================================================================
% GLOSSARY OF TERMS — Alert Aggregation in Microservices
% Standalone file for \input{glossary} in main IEEEtran document
% Compiles with: pdflatex main.tex (where main.tex includes this file)
% =============================================================================

% -----------------------------------------------------------------------------
% USAGE IN MAIN DOCUMENT:
% \section*{Glossary of Terms}
% \addcontentsline{toc}{section}{Glossary of Terms}
% \input{glossary}
% -----------------------------------------------------------------------------

\begin{description}

% =============================================================================
% CORE PROBLEM TERMS
% =============================================================================

\item[Alert Fatigue]
A cognitive state experienced by on-call engineers when the volume of alerts exceeds human cognitive capacity, leading to desensitization, missed critical alerts, and delayed incident response~\cite{pagerduty2020, grafana2025}.

\item[Alert Storm]
A sudden surge of alerts triggered by a single root cause propagating through dependent microservices, resulting in hundreds to thousands of alerts within minutes~\cite{kuang2024cola}.

\item[Noisy Alert Warning]
A warning-level alert generated by the monitoring system that is technically correct but does not require meaningful remediation action from an on-call engineer, classified into four categories: \textit{Duplicate}, \textit{Transient}, \textit{Unactionable}, and \textit{Cascading} (defined below).

\item[Early Warning Phase]
The pre-failure period during which a microservice exhibits performance degradation (e.g., increased latency, resource saturation) but has not yet violated Service Level Objectives (SLOs) or caused user-facing outages.

\item[Actionable Alert]
An alert that satisfies the Google SRE principle: ``every alert sent to an on-call engineer must require immediate human intervention''~\cite{google-sre-ch10, google-sre-ch11}.

% =============================================================================
% FOUR CATEGORIES OF NOISY ALERT WARNINGS
% =============================================================================

\item[Duplicate Warning Alert]
Multiple alerts originating from the same service with identical or near-identical message templates within a short time window ($\Delta t \le 60$~seconds), representing repeated observations of the same underlying condition rather than distinct incidents.

\item[Transient Warning Alert]
A warning alert that occurs during normal system operation (i.e., before fault injection time $T_{\text{inject}}$ in RCAEval) and resolves spontaneously without human intervention, typically caused by temporary network fluctuations or brief resource spikes.

\item[Unactionable Warning Alert]
A technically accurate warning alert for which no Standard Operating Procedure (SOP) or runbook exists, leaving the on-call engineer without a defined remediation path; the alert consumes attention but yields no operational value.

\item[Cascading Warning Alert]
A warning alert emitted by a downstream service that is a symptom of an upstream root cause; remediating the alerting service does not resolve the underlying issue, which originates in a different service (the root cause service).

% =============================================================================
% METHODOLOGY: FOUR EXPERIMENTAL ARMS
% =============================================================================

\item[Arm 1: Rule-Based Aggregation (Baseline)]
An alert grouping approach that clusters alerts based on static label matching (e.g., service name, alert name, severity) and fixed time windows (\texttt{group\_wait}, \texttt{group\_interval}), as implemented in Prometheus Alertmanager~\cite{prometheus-alertmanager}. Lacks awareness of service dependencies and semantic content.

\item[Arm 2: Semantic Similarity-Based Aggregation]
A method that parses raw logs into templates (via Drain3~\cite{drain3}), encodes templates into dense vector embeddings using Sentence-BERT~\cite{sbert}, and clusters alerts using density-based algorithms (HDBSCAN/DBSCAN) based on cosine similarity. Operates without topology knowledge.

\item[Arm 3: Temporal-Spatial Correlation Aggregation]
A topology-aware method that constructs a static service dependency graph (caller-callee adjacency matrix) from microservice architecture, then correlates alerts occurring within a sliding time window ($\Delta t \in \{30, 60, 120\}$~seconds) if a dependency path of length $k \le 2$ exists between the alerting services.

\item[Arm 4: LLM-Based Agentic Aggregation (Proposed)]
A hybrid two-stage approach: (1) a fast filter removes obvious duplicates and performs time-window pre-grouping; (2) an LLM agent (GPT-4o-mini / DeepSeek-V3) reasons over the remaining clusters using Chain-of-Thought prompting and three tools: \texttt{get\_related\_alerts}, \texttt{get\_service\_dependency}, and optionally \texttt{get\_runbook}. Limited to a maximum of two tool calls per grouping decision.

% =============================================================================
% TECHNICAL CONCEPTS & ARCHITECTURE
% =============================================================================

\item[Service Topology (Service Dependency Graph)]
A directed graph $G = (V, E)$ where vertices $V$ represent microservices and edges $E$ represent synchronous (HTTP/gRPC) or asynchronous (message queue) call relationships. Used to distinguish root cause services from downstream symptom services.

\item[Semantic Similarity]
A measure of meaning-based resemblance between two log messages, computed as the cosine similarity between their dense vector embeddings (typically 384-dimensional from \texttt{sentence-transformers/all-MiniLM-L6-v2}). Captures paraphrases and synonyms that lexical matching misses.

\item[Runbook Reasoning]
The process by which an LLM agent consults a synthetic runbook (a structured troubleshooting guide for a specific fault type) to determine whether an alert is actionable, requires escalation, or can be safely aggregated as noise.

\item[Tool-Calling Agent]
An LLM agent augmented with a predefined set of functions (tools) it can invoke during reasoning: \texttt{get\_related\_alerts(service, time\_window)}, \texttt{get\_service\_dependency(service)}, and \texttt{get\_runbook(fault\_type)}. Tool outputs are fed back into the agent's context for subsequent reasoning steps.

\item[Chain-of-Thought (CoT) Prompting]
A prompting technique that instructs the LLM to generate intermediate reasoning steps before producing a final answer, improving accuracy on multi-step causal inference tasks~\cite{wei2022cot}.

\item[Fast Filter (Deduplication + Time-Window Pre-Grouping)]
A lightweight deterministic preprocessing stage that removes exact duplicate alerts and performs initial grouping by service and time window, reducing the alert volume passed to the LLM agent by an estimated 70--80\%.

\item[Synthetic Runbook]
A researcher-authored troubleshooting document for each fault type in the benchmark dataset (e.g., CPU stress, memory leak, network latency), containing symptoms, diagnosis steps, and remediation actions. Not a production SOP; used solely for experimental evaluation.

% =============================================================================
% EVALUATION METRICS
% =============================================================================

\item[Alert Reduction Ratio (ARR)]
The primary metric quantifying alert volume reduction:
\begin{equation}
\text{ARR} = \left(1 - \frac{N_{\text{groups}}}{N_{\text{raw\_alerts}}}\right) \times 100\%
\end{equation}
where $N_{\text{raw\_alerts}}$ is the number of warning alerts before aggregation and $N_{\text{groups}}$ is the number of aggregated groups presented to the engineer. Target: $\text{ARR} \ge 60\%$.

\item[Root Cause Preservation Rate (RCPR)]
The critical safety metric ensuring the root cause evidence is not lost during aggregation. It is defined at two granularities. At service granularity, which is the primary metric and spans every case that exposes logs ($M = 359$):
\begin{equation}
\text{RCPR}_{\text{svc}} = \frac{1}{M} \sum_{i=1}^{M} \mathbb{I}(s_{\text{root}}^{(i)} \in \text{PrimaryGroup}_i) \times 100\%
\end{equation}
At alert granularity, a secondary metric restricted to the $M'$ cases whose ground-truth \texttt{root\_cause.txt} indicator line is itself a warning-tier log ($M' = 4$):
\begin{equation}
\text{RCPR}_{\text{alert}} = \frac{1}{M'} \sum_{i=1}^{M'} \mathbb{I}(\text{RootCauseAlert}_i \in \text{PrimaryGroup}_i) \times 100\%
\end{equation}
where $\mathbb{I}$ is the indicator function. Target: $\text{RCPR} \ge 95\%$ (mandatory).

\item[Pairwise Grouping Precision ($P_p$)]
The fraction of alert pairs placed in the same group that truly share the same root cause:
\begin{equation}
P_p = \frac{|\{(a_x, a_y) : \text{same group} \land \text{same root cause}\}|}{|\{(a_x, a_y) : \text{same group}\}|}
\end{equation}

\item[Pairwise Grouping Recall ($R_p$)]
The fraction of alert pairs sharing the same root cause that are successfully grouped together:
\begin{equation}
R_p = \frac{|\{(a_x, a_y) : \text{same group} \land \text{same root cause}\}|}{|\{(a_x, a_y) : \text{same root cause}\}|}
\end{equation}

\item[Pairwise Grouping F1-Score ($F1_p$)]
The harmonic mean of pairwise precision and recall:
\begin{equation}
F1_p = 2 \times \frac{P_p \times R_p}{P_p + R_p}
\end{equation}
Target: $F1_p \ge 0.85$.

\item[Group Purity]
The average homogeneity of root causes within each generated group:
\begin{equation}
\text{Purity} = \frac{1}{N_{\text{groups}}} \sum_{g=1}^{N_{\text{groups}}} \frac{\max_{c} |\{a \in g : \text{root\_cause}(a) = c\}|}{|g|}
\end{equation}
Target: $\text{Purity} \ge 0.80$.

\item[Processing Latency]
The wall-clock time required to process 100 warning alerts through the aggregation pipeline, measured in milliseconds. For Arm 4, includes LLM inference time and token overhead.

% =============================================================================
% DATASET & BENCHMARK TERMS
% =============================================================================

\item[RCAEval]
A benchmark for root cause analysis in microservice systems comprising 735 failure cases across three systems (Online Boutique, Sock Shop, Train Ticket) with logs, metrics, traces, and root cause labels~\cite{pham2025rcaeval}. Licensed under MIT. Direct inspection shows that only \textbf{359 of the 735 cases ship raw log telemetry} (suite RE1 is metric-only) and that the log schema contains \texttt{timestamp}, \texttt{container\_name}, and \texttt{message} but \textbf{no severity column}.

\item[Severity Tier Assignment]
The three-tier procedure by which this work assigns the severity tier $\ell_i$, since RCAEval supplies none: \emph{Tier 1 (explicit)} parses a level token from the message text (Sock Shop, Train Ticket); \emph{Tier 2 (frequency-inferred)} flags a template as anomalous when its fault-window frequency exceeds its normal-window frequency by a ratio $R$ (Online Boutique); \emph{Tier 3 (ground-truth marker)} tags the \texttt{root\_cause.txt} indicator line as \texttt{ROOT\_CAUSE} for the eight annotated cases.

\item[LEMMA-RCA]
A large-scale multi-modal dataset for root cause analysis containing hundreds of microservice entities and four fault types (DDoS, storage failure, resource contention, noisy neighbor)~\cite{zheng2024lemma}. Licensed under CC-BY-NC-4.0 (non-commercial). Its log pipeline aggregates raw log lines into three-dimensional time series using golden-signal keywords that cover only \texttt{error}, \texttt{exception}, and \texttt{critical} --- notably \textbf{not} \texttt{warning} --- so it offers no independent warning tier.

\item[Fault Injection]
The controlled introduction of a failure (e.g., CPU stress, network latency, packet loss) into a microservice system at a known timestamp $T_{\text{inject}}$ to generate labeled failure cases for benchmarking.

\item[Normal Window]
The interval $[T_{\text{start}}, T_{\text{inject}})$ preceding fault injection, during which the system operates normally; alerts occurring here are candidates for \textit{Transient} noise. Note that the dataset fields \texttt{normal\_timesteps} and \texttt{faulty\_timesteps} are integer counts of sample points, not timestamp lists, and therefore define these windows only indirectly.

\item[Fault Window]
The interval $[T_{\text{inject}}, T_{\text{end}}]$ following fault injection, during which the system exhibits degradation; alerts here are labeled as \textit{Signal} (if from the root cause service) or \textit{Cascading} (if from a service reachable from the root cause service in the call graph).

\item[Ground Truth Root Cause Label]
The authoritative annotation of which service is the root cause of a failure case, provided by the RCAEval benchmark. Used to compute RCPR and pairwise metrics.

% =============================================================================
% BASELINE & INFRASTRUCTURE TERMS
% =============================================================================

\item[Prometheus Alertmanager]
The de facto standard alert routing and grouping component in the Cloud Native Computing Foundation (CNCF) ecosystem. Groups alerts by label matchers and time windows (\texttt{group\_wait}, \texttt{group\_interval}, \texttt{group\_by})~\cite{prometheus-alertmanager}.

\item[Group Wait (\texttt{group\_wait})]
The initial delay after the first alert in a group before sending a notification, allowing additional alerts to arrive and be grouped. Default: 30~seconds.

\item[Group Interval (\texttt{group\_interval})]
The minimum interval between subsequent notifications for the same group. Default: 5~minutes.

\item[Drain3]
An online log parsing algorithm that extracts log templates from raw log streams by clustering log messages based on token similarity, producing parameterized templates with variable placeholders~\cite{drain3}.

\item[Sentence-BERT (S-BERT)]
A modification of BERT that uses siamese and triplet network structures to derive semantically meaningful sentence embeddings, optimized for cosine similarity comparison~\cite{sbert}.

\item[HDBSCAN]
Hierarchical Density-Based Spatial Clustering of Applications with Noise; a density-based clustering algorithm that does not require specifying the number of clusters and can identify outliers (noise points)~\cite{hdbscan}.

% =============================================================================
% EXPERIMENTAL DESIGN TERMS
% =============================================================================

\item[Ablation Study]
An experimental procedure where individual components (e.g., \texttt{get\_service\_dependency} tool, fast filter, runbook) are removed from the full system to measure their individual contribution to performance metrics.

\item[Cross-System Generalization]
The ability of an aggregation method trained or tuned on one microservice system (e.g., Online Boutique) to maintain performance on unseen systems (e.g., Train Ticket) without retraining.

\item[Statistical Significance Test]
Paired Wilcoxon signed-rank test (non-parametric) used to compare metric distributions across the four arms on the same failure cases, with significance level $\alpha = 0.05$.

\item[Metric Gaming]
The phenomenon where a method optimizes a single metric (e.g., ARR) at the expense of operational utility (e.g., grouping all alerts into one group achieves ARR $\approx$ 100\% but RCPR = 0\%). Mitigated by the multi-metric framework (ARR, RCPR, F1, Purity).

% =============================================================================
% INDUSTRY STANDARDS & REFERENCES
% =============================================================================

\item[Google SRE Principles]
Foundational Site Reliability Engineering practices documented in \textit{Site Reliability Engineering} (Ch.~6, 10, 11) and \textit{The Site Reliability Workbook} (Ch.~5), emphasizing actionable alerts, maximum 2 incidents per 12-hour shift, and alerting on SLOs rather than raw metrics~\cite{google-sre-book, google-sre-workbook}.

\item[Observability Pillars]
The three canonical telemetry types: \textit{Logs} (immutable event records), \textit{Metrics} (aggregated numerical measurements), and \textit{Traces} (distributed request flow records)~\cite{grafana2025}.

\end{description}

% =============================================================================
% LOCAL BIBLIOGRAPHY MACROS (if not using main document's .bib)
% Include these in your main .bib or define via IEEEabrv.bib/IEEEfull.bib
% =============================================================================
% @string{srebook = "Site Reliability Engineering: How Google Runs Production Systems"}
% @string{sreworkbook = "The Site Reliability Workbook: Practical Ways to Implement SRE"}
% @article{pagerduty2020, ...}
% @article{grafana2025, ...}
% @inproceedings{kuang2024cola, ...}
% @inproceedings{pham2025rcaeval, ...}
% @article{zheng2024lemma, ...}
% @inproceedings{yu2025alertguardian, ...}
% @manual{prometheus-alertmanager, ...}
% @article{sbert, ...}
% @article{drain3, ...}
% @article{hdbscan, ...}
% @article{wei2022cot, ...} in main IEEEtran document
% Compiles with: pdflatex main.tex (where main.tex includes this file)
% =============================================================================

% -----------------------------------------------------------------------------
% USAGE IN MAIN DOCUMENT:
% \section*{Glossary of Terms}
% \addcontentsline{toc}{section}{Glossary of Terms}
% % =============================================================================
% GLOSSARY OF TERMS — Alert Aggregation in Microservices
% Standalone file for \input{glossary} in main IEEEtran document
% Compiles with: pdflatex main.tex (where main.tex includes this file)
% =============================================================================

% -----------------------------------------------------------------------------
% USAGE IN MAIN DOCUMENT:
% \section*{Glossary of Terms}
% \addcontentsline{toc}{section}{Glossary of Terms}
% \input{glossary}
% -----------------------------------------------------------------------------

\begin{description}

% =============================================================================
% CORE PROBLEM TERMS
% =============================================================================

\item[Alert Fatigue]
A cognitive state experienced by on-call engineers when the volume of alerts exceeds human cognitive capacity, leading to desensitization, missed critical alerts, and delayed incident response~\cite{pagerduty2020, grafana2025}.

\item[Alert Storm]
A sudden surge of alerts triggered by a single root cause propagating through dependent microservices, resulting in hundreds to thousands of alerts within minutes~\cite{kuang2024cola}.

\item[Noisy Alert Warning]
A warning-level alert generated by the monitoring system that is technically correct but does not require meaningful remediation action from an on-call engineer, classified into four categories: \textit{Duplicate}, \textit{Transient}, \textit{Unactionable}, and \textit{Cascading} (defined below).

\item[Early Warning Phase]
The pre-failure period during which a microservice exhibits performance degradation (e.g., increased latency, resource saturation) but has not yet violated Service Level Objectives (SLOs) or caused user-facing outages.

\item[Actionable Alert]
An alert that satisfies the Google SRE principle: ``every alert sent to an on-call engineer must require immediate human intervention''~\cite{google-sre-ch10, google-sre-ch11}.

% =============================================================================
% FOUR CATEGORIES OF NOISY ALERT WARNINGS
% =============================================================================

\item[Duplicate Warning Alert]
Multiple alerts originating from the same service with identical or near-identical message templates within a short time window ($\Delta t \le 60$~seconds), representing repeated observations of the same underlying condition rather than distinct incidents.

\item[Transient Warning Alert]
A warning alert that occurs during normal system operation (i.e., before fault injection time $T_{\text{inject}}$ in RCAEval) and resolves spontaneously without human intervention, typically caused by temporary network fluctuations or brief resource spikes.

\item[Unactionable Warning Alert]
A technically accurate warning alert for which no Standard Operating Procedure (SOP) or runbook exists, leaving the on-call engineer without a defined remediation path; the alert consumes attention but yields no operational value.

\item[Cascading Warning Alert]
A warning alert emitted by a downstream service that is a symptom of an upstream root cause; remediating the alerting service does not resolve the underlying issue, which originates in a different service (the root cause service).

% =============================================================================
% METHODOLOGY: FOUR EXPERIMENTAL ARMS
% =============================================================================

\item[Arm 1: Rule-Based Aggregation (Baseline)]
An alert grouping approach that clusters alerts based on static label matching (e.g., service name, alert name, severity) and fixed time windows (\texttt{group\_wait}, \texttt{group\_interval}), as implemented in Prometheus Alertmanager~\cite{prometheus-alertmanager}. Lacks awareness of service dependencies and semantic content.

\item[Arm 2: Semantic Similarity-Based Aggregation]
A method that parses raw logs into templates (via Drain3~\cite{drain3}), encodes templates into dense vector embeddings using Sentence-BERT~\cite{sbert}, and clusters alerts using density-based algorithms (HDBSCAN/DBSCAN) based on cosine similarity. Operates without topology knowledge.

\item[Arm 3: Temporal-Spatial Correlation Aggregation]
A topology-aware method that constructs a static service dependency graph (caller-callee adjacency matrix) from microservice architecture, then correlates alerts occurring within a sliding time window ($\Delta t \in \{30, 60, 120\}$~seconds) if a dependency path of length $k \le 2$ exists between the alerting services.

\item[Arm 4: LLM-Based Agentic Aggregation (Proposed)]
A hybrid two-stage approach: (1) a fast filter removes obvious duplicates and performs time-window pre-grouping; (2) an LLM agent (GPT-4o-mini / DeepSeek-V3) reasons over the remaining clusters using Chain-of-Thought prompting and three tools: \texttt{get\_related\_alerts}, \texttt{get\_service\_dependency}, and optionally \texttt{get\_runbook}. Limited to a maximum of two tool calls per grouping decision.

% =============================================================================
% TECHNICAL CONCEPTS & ARCHITECTURE
% =============================================================================

\item[Service Topology (Service Dependency Graph)]
A directed graph $G = (V, E)$ where vertices $V$ represent microservices and edges $E$ represent synchronous (HTTP/gRPC) or asynchronous (message queue) call relationships. Used to distinguish root cause services from downstream symptom services.

\item[Semantic Similarity]
A measure of meaning-based resemblance between two log messages, computed as the cosine similarity between their dense vector embeddings (typically 384-dimensional from \texttt{sentence-transformers/all-MiniLM-L6-v2}). Captures paraphrases and synonyms that lexical matching misses.

\item[Runbook Reasoning]
The process by which an LLM agent consults a synthetic runbook (a structured troubleshooting guide for a specific fault type) to determine whether an alert is actionable, requires escalation, or can be safely aggregated as noise.

\item[Tool-Calling Agent]
An LLM agent augmented with a predefined set of functions (tools) it can invoke during reasoning: \texttt{get\_related\_alerts(service, time\_window)}, \texttt{get\_service\_dependency(service)}, and \texttt{get\_runbook(fault\_type)}. Tool outputs are fed back into the agent's context for subsequent reasoning steps.

\item[Chain-of-Thought (CoT) Prompting]
A prompting technique that instructs the LLM to generate intermediate reasoning steps before producing a final answer, improving accuracy on multi-step causal inference tasks~\cite{wei2022cot}.

\item[Fast Filter (Deduplication + Time-Window Pre-Grouping)]
A lightweight deterministic preprocessing stage that removes exact duplicate alerts and performs initial grouping by service and time window, reducing the alert volume passed to the LLM agent by an estimated 70--80\%.

\item[Synthetic Runbook]
A researcher-authored troubleshooting document for each fault type in the benchmark dataset (e.g., CPU stress, memory leak, network latency), containing symptoms, diagnosis steps, and remediation actions. Not a production SOP; used solely for experimental evaluation.

% =============================================================================
% EVALUATION METRICS
% =============================================================================

\item[Alert Reduction Ratio (ARR)]
The primary metric quantifying alert volume reduction:
\begin{equation}
\text{ARR} = \left(1 - \frac{N_{\text{groups}}}{N_{\text{raw\_alerts}}}\right) \times 100\%
\end{equation}
where $N_{\text{raw\_alerts}}$ is the number of warning alerts before aggregation and $N_{\text{groups}}$ is the number of aggregated groups presented to the engineer. Target: $\text{ARR} \ge 60\%$.

\item[Root Cause Preservation Rate (RCPR)]
The critical safety metric ensuring the root cause evidence is not lost during aggregation. It is defined at two granularities. At service granularity, which is the primary metric and spans every case that exposes logs ($M = 359$):
\begin{equation}
\text{RCPR}_{\text{svc}} = \frac{1}{M} \sum_{i=1}^{M} \mathbb{I}(s_{\text{root}}^{(i)} \in \text{PrimaryGroup}_i) \times 100\%
\end{equation}
At alert granularity, a secondary metric restricted to the $M'$ cases whose ground-truth \texttt{root\_cause.txt} indicator line is itself a warning-tier log ($M' = 4$):
\begin{equation}
\text{RCPR}_{\text{alert}} = \frac{1}{M'} \sum_{i=1}^{M'} \mathbb{I}(\text{RootCauseAlert}_i \in \text{PrimaryGroup}_i) \times 100\%
\end{equation}
where $\mathbb{I}$ is the indicator function. Target: $\text{RCPR} \ge 95\%$ (mandatory).

\item[Pairwise Grouping Precision ($P_p$)]
The fraction of alert pairs placed in the same group that truly share the same root cause:
\begin{equation}
P_p = \frac{|\{(a_x, a_y) : \text{same group} \land \text{same root cause}\}|}{|\{(a_x, a_y) : \text{same group}\}|}
\end{equation}

\item[Pairwise Grouping Recall ($R_p$)]
The fraction of alert pairs sharing the same root cause that are successfully grouped together:
\begin{equation}
R_p = \frac{|\{(a_x, a_y) : \text{same group} \land \text{same root cause}\}|}{|\{(a_x, a_y) : \text{same root cause}\}|}
\end{equation}

\item[Pairwise Grouping F1-Score ($F1_p$)]
The harmonic mean of pairwise precision and recall:
\begin{equation}
F1_p = 2 \times \frac{P_p \times R_p}{P_p + R_p}
\end{equation}
Target: $F1_p \ge 0.85$.

\item[Group Purity]
The average homogeneity of root causes within each generated group:
\begin{equation}
\text{Purity} = \frac{1}{N_{\text{groups}}} \sum_{g=1}^{N_{\text{groups}}} \frac{\max_{c} |\{a \in g : \text{root\_cause}(a) = c\}|}{|g|}
\end{equation}
Target: $\text{Purity} \ge 0.80$.

\item[Processing Latency]
The wall-clock time required to process 100 warning alerts through the aggregation pipeline, measured in milliseconds. For Arm 4, includes LLM inference time and token overhead.

% =============================================================================
% DATASET & BENCHMARK TERMS
% =============================================================================

\item[RCAEval]
A benchmark for root cause analysis in microservice systems comprising 735 failure cases across three systems (Online Boutique, Sock Shop, Train Ticket) with logs, metrics, traces, and root cause labels~\cite{pham2025rcaeval}. Licensed under MIT. Direct inspection shows that only \textbf{359 of the 735 cases ship raw log telemetry} (suite RE1 is metric-only) and that the log schema contains \texttt{timestamp}, \texttt{container\_name}, and \texttt{message} but \textbf{no severity column}.

\item[Severity Tier Assignment]
The three-tier procedure by which this work assigns the severity tier $\ell_i$, since RCAEval supplies none: \emph{Tier 1 (explicit)} parses a level token from the message text (Sock Shop, Train Ticket); \emph{Tier 2 (frequency-inferred)} flags a template as anomalous when its fault-window frequency exceeds its normal-window frequency by a ratio $R$ (Online Boutique); \emph{Tier 3 (ground-truth marker)} tags the \texttt{root\_cause.txt} indicator line as \texttt{ROOT\_CAUSE} for the eight annotated cases.

\item[LEMMA-RCA]
A large-scale multi-modal dataset for root cause analysis containing hundreds of microservice entities and four fault types (DDoS, storage failure, resource contention, noisy neighbor)~\cite{zheng2024lemma}. Licensed under CC-BY-NC-4.0 (non-commercial). Its log pipeline aggregates raw log lines into three-dimensional time series using golden-signal keywords that cover only \texttt{error}, \texttt{exception}, and \texttt{critical} --- notably \textbf{not} \texttt{warning} --- so it offers no independent warning tier.

\item[Fault Injection]
The controlled introduction of a failure (e.g., CPU stress, network latency, packet loss) into a microservice system at a known timestamp $T_{\text{inject}}$ to generate labeled failure cases for benchmarking.

\item[Normal Window]
The interval $[T_{\text{start}}, T_{\text{inject}})$ preceding fault injection, during which the system operates normally; alerts occurring here are candidates for \textit{Transient} noise. Note that the dataset fields \texttt{normal\_timesteps} and \texttt{faulty\_timesteps} are integer counts of sample points, not timestamp lists, and therefore define these windows only indirectly.

\item[Fault Window]
The interval $[T_{\text{inject}}, T_{\text{end}}]$ following fault injection, during which the system exhibits degradation; alerts here are labeled as \textit{Signal} (if from the root cause service) or \textit{Cascading} (if from a service reachable from the root cause service in the call graph).

\item[Ground Truth Root Cause Label]
The authoritative annotation of which service is the root cause of a failure case, provided by the RCAEval benchmark. Used to compute RCPR and pairwise metrics.

% =============================================================================
% BASELINE & INFRASTRUCTURE TERMS
% =============================================================================

\item[Prometheus Alertmanager]
The de facto standard alert routing and grouping component in the Cloud Native Computing Foundation (CNCF) ecosystem. Groups alerts by label matchers and time windows (\texttt{group\_wait}, \texttt{group\_interval}, \texttt{group\_by})~\cite{prometheus-alertmanager}.

\item[Group Wait (\texttt{group\_wait})]
The initial delay after the first alert in a group before sending a notification, allowing additional alerts to arrive and be grouped. Default: 30~seconds.

\item[Group Interval (\texttt{group\_interval})]
The minimum interval between subsequent notifications for the same group. Default: 5~minutes.

\item[Drain3]
An online log parsing algorithm that extracts log templates from raw log streams by clustering log messages based on token similarity, producing parameterized templates with variable placeholders~\cite{drain3}.

\item[Sentence-BERT (S-BERT)]
A modification of BERT that uses siamese and triplet network structures to derive semantically meaningful sentence embeddings, optimized for cosine similarity comparison~\cite{sbert}.

\item[HDBSCAN]
Hierarchical Density-Based Spatial Clustering of Applications with Noise; a density-based clustering algorithm that does not require specifying the number of clusters and can identify outliers (noise points)~\cite{hdbscan}.

% =============================================================================
% EXPERIMENTAL DESIGN TERMS
% =============================================================================

\item[Ablation Study]
An experimental procedure where individual components (e.g., \texttt{get\_service\_dependency} tool, fast filter, runbook) are removed from the full system to measure their individual contribution to performance metrics.

\item[Cross-System Generalization]
The ability of an aggregation method trained or tuned on one microservice system (e.g., Online Boutique) to maintain performance on unseen systems (e.g., Train Ticket) without retraining.

\item[Statistical Significance Test]
Paired Wilcoxon signed-rank test (non-parametric) used to compare metric distributions across the four arms on the same failure cases, with significance level $\alpha = 0.05$.

\item[Metric Gaming]
The phenomenon where a method optimizes a single metric (e.g., ARR) at the expense of operational utility (e.g., grouping all alerts into one group achieves ARR $\approx$ 100\% but RCPR = 0\%). Mitigated by the multi-metric framework (ARR, RCPR, F1, Purity).

% =============================================================================
% INDUSTRY STANDARDS & REFERENCES
% =============================================================================

\item[Google SRE Principles]
Foundational Site Reliability Engineering practices documented in \textit{Site Reliability Engineering} (Ch.~6, 10, 11) and \textit{The Site Reliability Workbook} (Ch.~5), emphasizing actionable alerts, maximum 2 incidents per 12-hour shift, and alerting on SLOs rather than raw metrics~\cite{google-sre-book, google-sre-workbook}.

\item[Observability Pillars]
The three canonical telemetry types: \textit{Logs} (immutable event records), \textit{Metrics} (aggregated numerical measurements), and \textit{Traces} (distributed request flow records)~\cite{grafana2025}.

\end{description}

% =============================================================================
% LOCAL BIBLIOGRAPHY MACROS (if not using main document's .bib)
% Include these in your main .bib or define via IEEEabrv.bib/IEEEfull.bib
% =============================================================================
% @string{srebook = "Site Reliability Engineering: How Google Runs Production Systems"}
% @string{sreworkbook = "The Site Reliability Workbook: Practical Ways to Implement SRE"}
% @article{pagerduty2020, ...}
% @article{grafana2025, ...}
% @inproceedings{kuang2024cola, ...}
% @inproceedings{pham2025rcaeval, ...}
% @article{zheng2024lemma, ...}
% @inproceedings{yu2025alertguardian, ...}
% @manual{prometheus-alertmanager, ...}
% @article{sbert, ...}
% @article{drain3, ...}
% @article{hdbscan, ...}
% @article{wei2022cot, ...}
% -----------------------------------------------------------------------------

\begin{description}

% =============================================================================
% CORE PROBLEM TERMS
% =============================================================================

\item[Alert Fatigue]
A cognitive state experienced by on-call engineers when the volume of alerts exceeds human cognitive capacity, leading to desensitization, missed critical alerts, and delayed incident response~\cite{pagerduty2020, grafana2025}.

\item[Alert Storm]
A sudden surge of alerts triggered by a single root cause propagating through dependent microservices, resulting in hundreds to thousands of alerts within minutes~\cite{kuang2024cola}.

\item[Noisy Alert Warning]
A warning-level alert generated by the monitoring system that is technically correct but does not require meaningful remediation action from an on-call engineer, classified into four categories: \textit{Duplicate}, \textit{Transient}, \textit{Unactionable}, and \textit{Cascading} (defined below).

\item[Early Warning Phase]
The pre-failure period during which a microservice exhibits performance degradation (e.g., increased latency, resource saturation) but has not yet violated Service Level Objectives (SLOs) or caused user-facing outages.

\item[Actionable Alert]
An alert that satisfies the Google SRE principle: ``every alert sent to an on-call engineer must require immediate human intervention''~\cite{google-sre-ch10, google-sre-ch11}.

% =============================================================================
% FOUR CATEGORIES OF NOISY ALERT WARNINGS
% =============================================================================

\item[Duplicate Warning Alert]
Multiple alerts originating from the same service with identical or near-identical message templates within a short time window ($\Delta t \le 60$~seconds), representing repeated observations of the same underlying condition rather than distinct incidents.

\item[Transient Warning Alert]
A warning alert that occurs during normal system operation (i.e., before fault injection time $T_{\text{inject}}$ in RCAEval) and resolves spontaneously without human intervention, typically caused by temporary network fluctuations or brief resource spikes.

\item[Unactionable Warning Alert]
A technically accurate warning alert for which no Standard Operating Procedure (SOP) or runbook exists, leaving the on-call engineer without a defined remediation path; the alert consumes attention but yields no operational value.

\item[Cascading Warning Alert]
A warning alert emitted by a downstream service that is a symptom of an upstream root cause; remediating the alerting service does not resolve the underlying issue, which originates in a different service (the root cause service).

% =============================================================================
% METHODOLOGY: FOUR EXPERIMENTAL ARMS
% =============================================================================

\item[Arm 1: Rule-Based Aggregation (Baseline)]
An alert grouping approach that clusters alerts based on static label matching (e.g., service name, alert name, severity) and fixed time windows (\texttt{group\_wait}, \texttt{group\_interval}), as implemented in Prometheus Alertmanager~\cite{prometheus-alertmanager}. Lacks awareness of service dependencies and semantic content.

\item[Arm 2: Semantic Similarity-Based Aggregation]
A method that parses raw logs into templates (via Drain3~\cite{drain3}), encodes templates into dense vector embeddings using Sentence-BERT~\cite{sbert}, and clusters alerts using density-based algorithms (HDBSCAN/DBSCAN) based on cosine similarity. Operates without topology knowledge.

\item[Arm 3: Temporal-Spatial Correlation Aggregation]
A topology-aware method that constructs a static service dependency graph (caller-callee adjacency matrix) from microservice architecture, then correlates alerts occurring within a sliding time window ($\Delta t \in \{30, 60, 120\}$~seconds) if a dependency path of length $k \le 2$ exists between the alerting services.

\item[Arm 4: LLM-Based Agentic Aggregation (Proposed)]
A hybrid two-stage approach: (1) a fast filter removes obvious duplicates and performs time-window pre-grouping; (2) an LLM agent (GPT-4o-mini / DeepSeek-V3) reasons over the remaining clusters using Chain-of-Thought prompting and three tools: \texttt{get\_related\_alerts}, \texttt{get\_service\_dependency}, and optionally \texttt{get\_runbook}. Limited to a maximum of two tool calls per grouping decision.

% =============================================================================
% TECHNICAL CONCEPTS & ARCHITECTURE
% =============================================================================

\item[Service Topology (Service Dependency Graph)]
A directed graph $G = (V, E)$ where vertices $V$ represent microservices and edges $E$ represent synchronous (HTTP/gRPC) or asynchronous (message queue) call relationships. Used to distinguish root cause services from downstream symptom services.

\item[Semantic Similarity]
A measure of meaning-based resemblance between two log messages, computed as the cosine similarity between their dense vector embeddings (typically 384-dimensional from \texttt{sentence-transformers/all-MiniLM-L6-v2}). Captures paraphrases and synonyms that lexical matching misses.

\item[Runbook Reasoning]
The process by which an LLM agent consults a synthetic runbook (a structured troubleshooting guide for a specific fault type) to determine whether an alert is actionable, requires escalation, or can be safely aggregated as noise.

\item[Tool-Calling Agent]
An LLM agent augmented with a predefined set of functions (tools) it can invoke during reasoning: \texttt{get\_related\_alerts(service, time\_window)}, \texttt{get\_service\_dependency(service)}, and \texttt{get\_runbook(fault\_type)}. Tool outputs are fed back into the agent's context for subsequent reasoning steps.

\item[Chain-of-Thought (CoT) Prompting]
A prompting technique that instructs the LLM to generate intermediate reasoning steps before producing a final answer, improving accuracy on multi-step causal inference tasks~\cite{wei2022cot}.

\item[Fast Filter (Deduplication + Time-Window Pre-Grouping)]
A lightweight deterministic preprocessing stage that removes exact duplicate alerts and performs initial grouping by service and time window, reducing the alert volume passed to the LLM agent by an estimated 70--80\%.

\item[Synthetic Runbook]
A researcher-authored troubleshooting document for each fault type in the benchmark dataset (e.g., CPU stress, memory leak, network latency), containing symptoms, diagnosis steps, and remediation actions. Not a production SOP; used solely for experimental evaluation.

% =============================================================================
% EVALUATION METRICS
% =============================================================================

\item[Alert Reduction Ratio (ARR)]
The primary metric quantifying alert volume reduction:
\begin{equation}
\text{ARR} = \left(1 - \frac{N_{\text{groups}}}{N_{\text{raw\_alerts}}}\right) \times 100\%
\end{equation}
where $N_{\text{raw\_alerts}}$ is the number of warning alerts before aggregation and $N_{\text{groups}}$ is the number of aggregated groups presented to the engineer. Target: $\text{ARR} \ge 60\%$.

\item[Root Cause Preservation Rate (RCPR)]
The critical safety metric ensuring the root cause evidence is not lost during aggregation. It is defined at two granularities. At service granularity, which is the primary metric and spans every case that exposes logs ($M = 359$):
\begin{equation}
\text{RCPR}_{\text{svc}} = \frac{1}{M} \sum_{i=1}^{M} \mathbb{I}(s_{\text{root}}^{(i)} \in \text{PrimaryGroup}_i) \times 100\%
\end{equation}
At alert granularity, a secondary metric restricted to the $M'$ cases whose ground-truth \texttt{root\_cause.txt} indicator line is itself a warning-tier log ($M' = 4$):
\begin{equation}
\text{RCPR}_{\text{alert}} = \frac{1}{M'} \sum_{i=1}^{M'} \mathbb{I}(\text{RootCauseAlert}_i \in \text{PrimaryGroup}_i) \times 100\%
\end{equation}
where $\mathbb{I}$ is the indicator function. Target: $\text{RCPR} \ge 95\%$ (mandatory).

\item[Pairwise Grouping Precision ($P_p$)]
The fraction of alert pairs placed in the same group that truly share the same root cause:
\begin{equation}
P_p = \frac{|\{(a_x, a_y) : \text{same group} \land \text{same root cause}\}|}{|\{(a_x, a_y) : \text{same group}\}|}
\end{equation}

\item[Pairwise Grouping Recall ($R_p$)]
The fraction of alert pairs sharing the same root cause that are successfully grouped together:
\begin{equation}
R_p = \frac{|\{(a_x, a_y) : \text{same group} \land \text{same root cause}\}|}{|\{(a_x, a_y) : \text{same root cause}\}|}
\end{equation}

\item[Pairwise Grouping F1-Score ($F1_p$)]
The harmonic mean of pairwise precision and recall:
\begin{equation}
F1_p = 2 \times \frac{P_p \times R_p}{P_p + R_p}
\end{equation}
Target: $F1_p \ge 0.85$.

\item[Group Purity]
The average homogeneity of root causes within each generated group:
\begin{equation}
\text{Purity} = \frac{1}{N_{\text{groups}}} \sum_{g=1}^{N_{\text{groups}}} \frac{\max_{c} |\{a \in g : \text{root\_cause}(a) = c\}|}{|g|}
\end{equation}
Target: $\text{Purity} \ge 0.80$.

\item[Processing Latency]
The wall-clock time required to process 100 warning alerts through the aggregation pipeline, measured in milliseconds. For Arm 4, includes LLM inference time and token overhead.

% =============================================================================
% DATASET & BENCHMARK TERMS
% =============================================================================

\item[RCAEval]
A benchmark for root cause analysis in microservice systems comprising 735 failure cases across three systems (Online Boutique, Sock Shop, Train Ticket) with logs, metrics, traces, and root cause labels~\cite{pham2025rcaeval}. Licensed under MIT. Direct inspection shows that only \textbf{359 of the 735 cases ship raw log telemetry} (suite RE1 is metric-only) and that the log schema contains \texttt{timestamp}, \texttt{container\_name}, and \texttt{message} but \textbf{no severity column}.

\item[Severity Tier Assignment]
The three-tier procedure by which this work assigns the severity tier $\ell_i$, since RCAEval supplies none: \emph{Tier 1 (explicit)} parses a level token from the message text (Sock Shop, Train Ticket); \emph{Tier 2 (frequency-inferred)} flags a template as anomalous when its fault-window frequency exceeds its normal-window frequency by a ratio $R$ (Online Boutique); \emph{Tier 3 (ground-truth marker)} tags the \texttt{root\_cause.txt} indicator line as \texttt{ROOT\_CAUSE} for the eight annotated cases.

\item[LEMMA-RCA]
A large-scale multi-modal dataset for root cause analysis containing hundreds of microservice entities and four fault types (DDoS, storage failure, resource contention, noisy neighbor)~\cite{zheng2024lemma}. Licensed under CC-BY-NC-4.0 (non-commercial). Its log pipeline aggregates raw log lines into three-dimensional time series using golden-signal keywords that cover only \texttt{error}, \texttt{exception}, and \texttt{critical} --- notably \textbf{not} \texttt{warning} --- so it offers no independent warning tier.

\item[Fault Injection]
The controlled introduction of a failure (e.g., CPU stress, network latency, packet loss) into a microservice system at a known timestamp $T_{\text{inject}}$ to generate labeled failure cases for benchmarking.

\item[Normal Window]
The interval $[T_{\text{start}}, T_{\text{inject}})$ preceding fault injection, during which the system operates normally; alerts occurring here are candidates for \textit{Transient} noise. Note that the dataset fields \texttt{normal\_timesteps} and \texttt{faulty\_timesteps} are integer counts of sample points, not timestamp lists, and therefore define these windows only indirectly.

\item[Fault Window]
The interval $[T_{\text{inject}}, T_{\text{end}}]$ following fault injection, during which the system exhibits degradation; alerts here are labeled as \textit{Signal} (if from the root cause service) or \textit{Cascading} (if from a service reachable from the root cause service in the call graph).

\item[Ground Truth Root Cause Label]
The authoritative annotation of which service is the root cause of a failure case, provided by the RCAEval benchmark. Used to compute RCPR and pairwise metrics.

% =============================================================================
% BASELINE & INFRASTRUCTURE TERMS
% =============================================================================

\item[Prometheus Alertmanager]
The de facto standard alert routing and grouping component in the Cloud Native Computing Foundation (CNCF) ecosystem. Groups alerts by label matchers and time windows (\texttt{group\_wait}, \texttt{group\_interval}, \texttt{group\_by})~\cite{prometheus-alertmanager}.

\item[Group Wait (\texttt{group\_wait})]
The initial delay after the first alert in a group before sending a notification, allowing additional alerts to arrive and be grouped. Default: 30~seconds.

\item[Group Interval (\texttt{group\_interval})]
The minimum interval between subsequent notifications for the same group. Default: 5~minutes.

\item[Drain3]
An online log parsing algorithm that extracts log templates from raw log streams by clustering log messages based on token similarity, producing parameterized templates with variable placeholders~\cite{drain3}.

\item[Sentence-BERT (S-BERT)]
A modification of BERT that uses siamese and triplet network structures to derive semantically meaningful sentence embeddings, optimized for cosine similarity comparison~\cite{sbert}.

\item[HDBSCAN]
Hierarchical Density-Based Spatial Clustering of Applications with Noise; a density-based clustering algorithm that does not require specifying the number of clusters and can identify outliers (noise points)~\cite{hdbscan}.

% =============================================================================
% EXPERIMENTAL DESIGN TERMS
% =============================================================================

\item[Ablation Study]
An experimental procedure where individual components (e.g., \texttt{get\_service\_dependency} tool, fast filter, runbook) are removed from the full system to measure their individual contribution to performance metrics.

\item[Cross-System Generalization]
The ability of an aggregation method trained or tuned on one microservice system (e.g., Online Boutique) to maintain performance on unseen systems (e.g., Train Ticket) without retraining.

\item[Statistical Significance Test]
Paired Wilcoxon signed-rank test (non-parametric) used to compare metric distributions across the four arms on the same failure cases, with significance level $\alpha = 0.05$.

\item[Metric Gaming]
The phenomenon where a method optimizes a single metric (e.g., ARR) at the expense of operational utility (e.g., grouping all alerts into one group achieves ARR $\approx$ 100\% but RCPR = 0\%). Mitigated by the multi-metric framework (ARR, RCPR, F1, Purity).

% =============================================================================
% INDUSTRY STANDARDS & REFERENCES
% =============================================================================

\item[Google SRE Principles]
Foundational Site Reliability Engineering practices documented in \textit{Site Reliability Engineering} (Ch.~6, 10, 11) and \textit{The Site Reliability Workbook} (Ch.~5), emphasizing actionable alerts, maximum 2 incidents per 12-hour shift, and alerting on SLOs rather than raw metrics~\cite{google-sre-book, google-sre-workbook}.

\item[Observability Pillars]
The three canonical telemetry types: \textit{Logs} (immutable event records), \textit{Metrics} (aggregated numerical measurements), and \textit{Traces} (distributed request flow records)~\cite{grafana2025}.

\end{description}

% =============================================================================
% LOCAL BIBLIOGRAPHY MACROS (if not using main document's .bib)
% Include these in your main .bib or define via IEEEabrv.bib/IEEEfull.bib
% =============================================================================
% @string{srebook = "Site Reliability Engineering: How Google Runs Production Systems"}
% @string{sreworkbook = "The Site Reliability Workbook: Practical Ways to Implement SRE"}
% @article{pagerduty2020, ...}
% @article{grafana2025, ...}
% @inproceedings{kuang2024cola, ...}
% @inproceedings{pham2025rcaeval, ...}
% @article{zheng2024lemma, ...}
% @inproceedings{yu2025alertguardian, ...}
% @manual{prometheus-alertmanager, ...}
% @article{sbert, ...}
% @article{drain3, ...}
% @article{hdbscan, ...}
% @article{wei2022cot, ...} in main IEEEtran document
% Compiles with: pdflatex main.tex (where main.tex includes this file)
% =============================================================================

% -----------------------------------------------------------------------------
% USAGE IN MAIN DOCUMENT:
% \section*{Glossary of Terms}
% \addcontentsline{toc}{section}{Glossary of Terms}
% % =============================================================================
% GLOSSARY OF TERMS — Alert Aggregation in Microservices
% Standalone file for % =============================================================================
% GLOSSARY OF TERMS — Alert Aggregation in Microservices
% Standalone file for \input{glossary} in main IEEEtran document
% Compiles with: pdflatex main.tex (where main.tex includes this file)
% =============================================================================

% -----------------------------------------------------------------------------
% USAGE IN MAIN DOCUMENT:
% \section*{Glossary of Terms}
% \addcontentsline{toc}{section}{Glossary of Terms}
% \input{glossary}
% -----------------------------------------------------------------------------

\begin{description}

% =============================================================================
% CORE PROBLEM TERMS
% =============================================================================

\item[Alert Fatigue]
A cognitive state experienced by on-call engineers when the volume of alerts exceeds human cognitive capacity, leading to desensitization, missed critical alerts, and delayed incident response~\cite{pagerduty2020, grafana2025}.

\item[Alert Storm]
A sudden surge of alerts triggered by a single root cause propagating through dependent microservices, resulting in hundreds to thousands of alerts within minutes~\cite{kuang2024cola}.

\item[Noisy Alert Warning]
A warning-level alert generated by the monitoring system that is technically correct but does not require meaningful remediation action from an on-call engineer, classified into four categories: \textit{Duplicate}, \textit{Transient}, \textit{Unactionable}, and \textit{Cascading} (defined below).

\item[Early Warning Phase]
The pre-failure period during which a microservice exhibits performance degradation (e.g., increased latency, resource saturation) but has not yet violated Service Level Objectives (SLOs) or caused user-facing outages.

\item[Actionable Alert]
An alert that satisfies the Google SRE principle: ``every alert sent to an on-call engineer must require immediate human intervention''~\cite{google-sre-ch10, google-sre-ch11}.

% =============================================================================
% FOUR CATEGORIES OF NOISY ALERT WARNINGS
% =============================================================================

\item[Duplicate Warning Alert]
Multiple alerts originating from the same service with identical or near-identical message templates within a short time window ($\Delta t \le 60$~seconds), representing repeated observations of the same underlying condition rather than distinct incidents.

\item[Transient Warning Alert]
A warning alert that occurs during normal system operation (i.e., before fault injection time $T_{\text{inject}}$ in RCAEval) and resolves spontaneously without human intervention, typically caused by temporary network fluctuations or brief resource spikes.

\item[Unactionable Warning Alert]
A technically accurate warning alert for which no Standard Operating Procedure (SOP) or runbook exists, leaving the on-call engineer without a defined remediation path; the alert consumes attention but yields no operational value.

\item[Cascading Warning Alert]
A warning alert emitted by a downstream service that is a symptom of an upstream root cause; remediating the alerting service does not resolve the underlying issue, which originates in a different service (the root cause service).

% =============================================================================
% METHODOLOGY: FOUR EXPERIMENTAL ARMS
% =============================================================================

\item[Arm 1: Rule-Based Aggregation (Baseline)]
An alert grouping approach that clusters alerts based on static label matching (e.g., service name, alert name, severity) and fixed time windows (\texttt{group\_wait}, \texttt{group\_interval}), as implemented in Prometheus Alertmanager~\cite{prometheus-alertmanager}. Lacks awareness of service dependencies and semantic content.

\item[Arm 2: Semantic Similarity-Based Aggregation]
A method that parses raw logs into templates (via Drain3~\cite{drain3}), encodes templates into dense vector embeddings using Sentence-BERT~\cite{sbert}, and clusters alerts using density-based algorithms (HDBSCAN/DBSCAN) based on cosine similarity. Operates without topology knowledge.

\item[Arm 3: Temporal-Spatial Correlation Aggregation]
A topology-aware method that constructs a static service dependency graph (caller-callee adjacency matrix) from microservice architecture, then correlates alerts occurring within a sliding time window ($\Delta t \in \{30, 60, 120\}$~seconds) if a dependency path of length $k \le 2$ exists between the alerting services.

\item[Arm 4: LLM-Based Agentic Aggregation (Proposed)]
A hybrid two-stage approach: (1) a fast filter removes obvious duplicates and performs time-window pre-grouping; (2) an LLM agent (GPT-4o-mini / DeepSeek-V3) reasons over the remaining clusters using Chain-of-Thought prompting and three tools: \texttt{get\_related\_alerts}, \texttt{get\_service\_dependency}, and optionally \texttt{get\_runbook}. Limited to a maximum of two tool calls per grouping decision.

% =============================================================================
% TECHNICAL CONCEPTS & ARCHITECTURE
% =============================================================================

\item[Service Topology (Service Dependency Graph)]
A directed graph $G = (V, E)$ where vertices $V$ represent microservices and edges $E$ represent synchronous (HTTP/gRPC) or asynchronous (message queue) call relationships. Used to distinguish root cause services from downstream symptom services.

\item[Semantic Similarity]
A measure of meaning-based resemblance between two log messages, computed as the cosine similarity between their dense vector embeddings (typically 384-dimensional from \texttt{sentence-transformers/all-MiniLM-L6-v2}). Captures paraphrases and synonyms that lexical matching misses.

\item[Runbook Reasoning]
The process by which an LLM agent consults a synthetic runbook (a structured troubleshooting guide for a specific fault type) to determine whether an alert is actionable, requires escalation, or can be safely aggregated as noise.

\item[Tool-Calling Agent]
An LLM agent augmented with a predefined set of functions (tools) it can invoke during reasoning: \texttt{get\_related\_alerts(service, time\_window)}, \texttt{get\_service\_dependency(service)}, and \texttt{get\_runbook(fault\_type)}. Tool outputs are fed back into the agent's context for subsequent reasoning steps.

\item[Chain-of-Thought (CoT) Prompting]
A prompting technique that instructs the LLM to generate intermediate reasoning steps before producing a final answer, improving accuracy on multi-step causal inference tasks~\cite{wei2022cot}.

\item[Fast Filter (Deduplication + Time-Window Pre-Grouping)]
A lightweight deterministic preprocessing stage that removes exact duplicate alerts and performs initial grouping by service and time window, reducing the alert volume passed to the LLM agent by an estimated 70--80\%.

\item[Synthetic Runbook]
A researcher-authored troubleshooting document for each fault type in the benchmark dataset (e.g., CPU stress, memory leak, network latency), containing symptoms, diagnosis steps, and remediation actions. Not a production SOP; used solely for experimental evaluation.

% =============================================================================
% EVALUATION METRICS
% =============================================================================

\item[Alert Reduction Ratio (ARR)]
The primary metric quantifying alert volume reduction:
\begin{equation}
\text{ARR} = \left(1 - \frac{N_{\text{groups}}}{N_{\text{raw\_alerts}}}\right) \times 100\%
\end{equation}
where $N_{\text{raw\_alerts}}$ is the number of warning alerts before aggregation and $N_{\text{groups}}$ is the number of aggregated groups presented to the engineer. Target: $\text{ARR} \ge 60\%$.

\item[Root Cause Preservation Rate (RCPR)]
The critical safety metric ensuring the root cause evidence is not lost during aggregation. It is defined at two granularities. At service granularity, which is the primary metric and spans every case that exposes logs ($M = 359$):
\begin{equation}
\text{RCPR}_{\text{svc}} = \frac{1}{M} \sum_{i=1}^{M} \mathbb{I}(s_{\text{root}}^{(i)} \in \text{PrimaryGroup}_i) \times 100\%
\end{equation}
At alert granularity, a secondary metric restricted to the $M'$ cases whose ground-truth \texttt{root\_cause.txt} indicator line is itself a warning-tier log ($M' = 4$):
\begin{equation}
\text{RCPR}_{\text{alert}} = \frac{1}{M'} \sum_{i=1}^{M'} \mathbb{I}(\text{RootCauseAlert}_i \in \text{PrimaryGroup}_i) \times 100\%
\end{equation}
where $\mathbb{I}$ is the indicator function. Target: $\text{RCPR} \ge 95\%$ (mandatory).

\item[Pairwise Grouping Precision ($P_p$)]
The fraction of alert pairs placed in the same group that truly share the same root cause:
\begin{equation}
P_p = \frac{|\{(a_x, a_y) : \text{same group} \land \text{same root cause}\}|}{|\{(a_x, a_y) : \text{same group}\}|}
\end{equation}

\item[Pairwise Grouping Recall ($R_p$)]
The fraction of alert pairs sharing the same root cause that are successfully grouped together:
\begin{equation}
R_p = \frac{|\{(a_x, a_y) : \text{same group} \land \text{same root cause}\}|}{|\{(a_x, a_y) : \text{same root cause}\}|}
\end{equation}

\item[Pairwise Grouping F1-Score ($F1_p$)]
The harmonic mean of pairwise precision and recall:
\begin{equation}
F1_p = 2 \times \frac{P_p \times R_p}{P_p + R_p}
\end{equation}
Target: $F1_p \ge 0.85$.

\item[Group Purity]
The average homogeneity of root causes within each generated group:
\begin{equation}
\text{Purity} = \frac{1}{N_{\text{groups}}} \sum_{g=1}^{N_{\text{groups}}} \frac{\max_{c} |\{a \in g : \text{root\_cause}(a) = c\}|}{|g|}
\end{equation}
Target: $\text{Purity} \ge 0.80$.

\item[Processing Latency]
The wall-clock time required to process 100 warning alerts through the aggregation pipeline, measured in milliseconds. For Arm 4, includes LLM inference time and token overhead.

% =============================================================================
% DATASET & BENCHMARK TERMS
% =============================================================================

\item[RCAEval]
A benchmark for root cause analysis in microservice systems comprising 735 failure cases across three systems (Online Boutique, Sock Shop, Train Ticket) with logs, metrics, traces, and root cause labels~\cite{pham2025rcaeval}. Licensed under MIT. Direct inspection shows that only \textbf{359 of the 735 cases ship raw log telemetry} (suite RE1 is metric-only) and that the log schema contains \texttt{timestamp}, \texttt{container\_name}, and \texttt{message} but \textbf{no severity column}.

\item[Severity Tier Assignment]
The three-tier procedure by which this work assigns the severity tier $\ell_i$, since RCAEval supplies none: \emph{Tier 1 (explicit)} parses a level token from the message text (Sock Shop, Train Ticket); \emph{Tier 2 (frequency-inferred)} flags a template as anomalous when its fault-window frequency exceeds its normal-window frequency by a ratio $R$ (Online Boutique); \emph{Tier 3 (ground-truth marker)} tags the \texttt{root\_cause.txt} indicator line as \texttt{ROOT\_CAUSE} for the eight annotated cases.

\item[LEMMA-RCA]
A large-scale multi-modal dataset for root cause analysis containing hundreds of microservice entities and four fault types (DDoS, storage failure, resource contention, noisy neighbor)~\cite{zheng2024lemma}. Licensed under CC-BY-NC-4.0 (non-commercial). Its log pipeline aggregates raw log lines into three-dimensional time series using golden-signal keywords that cover only \texttt{error}, \texttt{exception}, and \texttt{critical} --- notably \textbf{not} \texttt{warning} --- so it offers no independent warning tier.

\item[Fault Injection]
The controlled introduction of a failure (e.g., CPU stress, network latency, packet loss) into a microservice system at a known timestamp $T_{\text{inject}}$ to generate labeled failure cases for benchmarking.

\item[Normal Window]
The interval $[T_{\text{start}}, T_{\text{inject}})$ preceding fault injection, during which the system operates normally; alerts occurring here are candidates for \textit{Transient} noise. Note that the dataset fields \texttt{normal\_timesteps} and \texttt{faulty\_timesteps} are integer counts of sample points, not timestamp lists, and therefore define these windows only indirectly.

\item[Fault Window]
The interval $[T_{\text{inject}}, T_{\text{end}}]$ following fault injection, during which the system exhibits degradation; alerts here are labeled as \textit{Signal} (if from the root cause service) or \textit{Cascading} (if from a service reachable from the root cause service in the call graph).

\item[Ground Truth Root Cause Label]
The authoritative annotation of which service is the root cause of a failure case, provided by the RCAEval benchmark. Used to compute RCPR and pairwise metrics.

% =============================================================================
% BASELINE & INFRASTRUCTURE TERMS
% =============================================================================

\item[Prometheus Alertmanager]
The de facto standard alert routing and grouping component in the Cloud Native Computing Foundation (CNCF) ecosystem. Groups alerts by label matchers and time windows (\texttt{group\_wait}, \texttt{group\_interval}, \texttt{group\_by})~\cite{prometheus-alertmanager}.

\item[Group Wait (\texttt{group\_wait})]
The initial delay after the first alert in a group before sending a notification, allowing additional alerts to arrive and be grouped. Default: 30~seconds.

\item[Group Interval (\texttt{group\_interval})]
The minimum interval between subsequent notifications for the same group. Default: 5~minutes.

\item[Drain3]
An online log parsing algorithm that extracts log templates from raw log streams by clustering log messages based on token similarity, producing parameterized templates with variable placeholders~\cite{drain3}.

\item[Sentence-BERT (S-BERT)]
A modification of BERT that uses siamese and triplet network structures to derive semantically meaningful sentence embeddings, optimized for cosine similarity comparison~\cite{sbert}.

\item[HDBSCAN]
Hierarchical Density-Based Spatial Clustering of Applications with Noise; a density-based clustering algorithm that does not require specifying the number of clusters and can identify outliers (noise points)~\cite{hdbscan}.

% =============================================================================
% EXPERIMENTAL DESIGN TERMS
% =============================================================================

\item[Ablation Study]
An experimental procedure where individual components (e.g., \texttt{get\_service\_dependency} tool, fast filter, runbook) are removed from the full system to measure their individual contribution to performance metrics.

\item[Cross-System Generalization]
The ability of an aggregation method trained or tuned on one microservice system (e.g., Online Boutique) to maintain performance on unseen systems (e.g., Train Ticket) without retraining.

\item[Statistical Significance Test]
Paired Wilcoxon signed-rank test (non-parametric) used to compare metric distributions across the four arms on the same failure cases, with significance level $\alpha = 0.05$.

\item[Metric Gaming]
The phenomenon where a method optimizes a single metric (e.g., ARR) at the expense of operational utility (e.g., grouping all alerts into one group achieves ARR $\approx$ 100\% but RCPR = 0\%). Mitigated by the multi-metric framework (ARR, RCPR, F1, Purity).

% =============================================================================
% INDUSTRY STANDARDS & REFERENCES
% =============================================================================

\item[Google SRE Principles]
Foundational Site Reliability Engineering practices documented in \textit{Site Reliability Engineering} (Ch.~6, 10, 11) and \textit{The Site Reliability Workbook} (Ch.~5), emphasizing actionable alerts, maximum 2 incidents per 12-hour shift, and alerting on SLOs rather than raw metrics~\cite{google-sre-book, google-sre-workbook}.

\item[Observability Pillars]
The three canonical telemetry types: \textit{Logs} (immutable event records), \textit{Metrics} (aggregated numerical measurements), and \textit{Traces} (distributed request flow records)~\cite{grafana2025}.

\end{description}

% =============================================================================
% LOCAL BIBLIOGRAPHY MACROS (if not using main document's .bib)
% Include these in your main .bib or define via IEEEabrv.bib/IEEEfull.bib
% =============================================================================
% @string{srebook = "Site Reliability Engineering: How Google Runs Production Systems"}
% @string{sreworkbook = "The Site Reliability Workbook: Practical Ways to Implement SRE"}
% @article{pagerduty2020, ...}
% @article{grafana2025, ...}
% @inproceedings{kuang2024cola, ...}
% @inproceedings{pham2025rcaeval, ...}
% @article{zheng2024lemma, ...}
% @inproceedings{yu2025alertguardian, ...}
% @manual{prometheus-alertmanager, ...}
% @article{sbert, ...}
% @article{drain3, ...}
% @article{hdbscan, ...}
% @article{wei2022cot, ...} in main IEEEtran document
% Compiles with: pdflatex main.tex (where main.tex includes this file)
% =============================================================================

% -----------------------------------------------------------------------------
% USAGE IN MAIN DOCUMENT:
% \section*{Glossary of Terms}
% \addcontentsline{toc}{section}{Glossary of Terms}
% % =============================================================================
% GLOSSARY OF TERMS — Alert Aggregation in Microservices
% Standalone file for \input{glossary} in main IEEEtran document
% Compiles with: pdflatex main.tex (where main.tex includes this file)
% =============================================================================

% -----------------------------------------------------------------------------
% USAGE IN MAIN DOCUMENT:
% \section*{Glossary of Terms}
% \addcontentsline{toc}{section}{Glossary of Terms}
% \input{glossary}
% -----------------------------------------------------------------------------

\begin{description}

% =============================================================================
% CORE PROBLEM TERMS
% =============================================================================

\item[Alert Fatigue]
A cognitive state experienced by on-call engineers when the volume of alerts exceeds human cognitive capacity, leading to desensitization, missed critical alerts, and delayed incident response~\cite{pagerduty2020, grafana2025}.

\item[Alert Storm]
A sudden surge of alerts triggered by a single root cause propagating through dependent microservices, resulting in hundreds to thousands of alerts within minutes~\cite{kuang2024cola}.

\item[Noisy Alert Warning]
A warning-level alert generated by the monitoring system that is technically correct but does not require meaningful remediation action from an on-call engineer, classified into four categories: \textit{Duplicate}, \textit{Transient}, \textit{Unactionable}, and \textit{Cascading} (defined below).

\item[Early Warning Phase]
The pre-failure period during which a microservice exhibits performance degradation (e.g., increased latency, resource saturation) but has not yet violated Service Level Objectives (SLOs) or caused user-facing outages.

\item[Actionable Alert]
An alert that satisfies the Google SRE principle: ``every alert sent to an on-call engineer must require immediate human intervention''~\cite{google-sre-ch10, google-sre-ch11}.

% =============================================================================
% FOUR CATEGORIES OF NOISY ALERT WARNINGS
% =============================================================================

\item[Duplicate Warning Alert]
Multiple alerts originating from the same service with identical or near-identical message templates within a short time window ($\Delta t \le 60$~seconds), representing repeated observations of the same underlying condition rather than distinct incidents.

\item[Transient Warning Alert]
A warning alert that occurs during normal system operation (i.e., before fault injection time $T_{\text{inject}}$ in RCAEval) and resolves spontaneously without human intervention, typically caused by temporary network fluctuations or brief resource spikes.

\item[Unactionable Warning Alert]
A technically accurate warning alert for which no Standard Operating Procedure (SOP) or runbook exists, leaving the on-call engineer without a defined remediation path; the alert consumes attention but yields no operational value.

\item[Cascading Warning Alert]
A warning alert emitted by a downstream service that is a symptom of an upstream root cause; remediating the alerting service does not resolve the underlying issue, which originates in a different service (the root cause service).

% =============================================================================
% METHODOLOGY: FOUR EXPERIMENTAL ARMS
% =============================================================================

\item[Arm 1: Rule-Based Aggregation (Baseline)]
An alert grouping approach that clusters alerts based on static label matching (e.g., service name, alert name, severity) and fixed time windows (\texttt{group\_wait}, \texttt{group\_interval}), as implemented in Prometheus Alertmanager~\cite{prometheus-alertmanager}. Lacks awareness of service dependencies and semantic content.

\item[Arm 2: Semantic Similarity-Based Aggregation]
A method that parses raw logs into templates (via Drain3~\cite{drain3}), encodes templates into dense vector embeddings using Sentence-BERT~\cite{sbert}, and clusters alerts using density-based algorithms (HDBSCAN/DBSCAN) based on cosine similarity. Operates without topology knowledge.

\item[Arm 3: Temporal-Spatial Correlation Aggregation]
A topology-aware method that constructs a static service dependency graph (caller-callee adjacency matrix) from microservice architecture, then correlates alerts occurring within a sliding time window ($\Delta t \in \{30, 60, 120\}$~seconds) if a dependency path of length $k \le 2$ exists between the alerting services.

\item[Arm 4: LLM-Based Agentic Aggregation (Proposed)]
A hybrid two-stage approach: (1) a fast filter removes obvious duplicates and performs time-window pre-grouping; (2) an LLM agent (GPT-4o-mini / DeepSeek-V3) reasons over the remaining clusters using Chain-of-Thought prompting and three tools: \texttt{get\_related\_alerts}, \texttt{get\_service\_dependency}, and optionally \texttt{get\_runbook}. Limited to a maximum of two tool calls per grouping decision.

% =============================================================================
% TECHNICAL CONCEPTS & ARCHITECTURE
% =============================================================================

\item[Service Topology (Service Dependency Graph)]
A directed graph $G = (V, E)$ where vertices $V$ represent microservices and edges $E$ represent synchronous (HTTP/gRPC) or asynchronous (message queue) call relationships. Used to distinguish root cause services from downstream symptom services.

\item[Semantic Similarity]
A measure of meaning-based resemblance between two log messages, computed as the cosine similarity between their dense vector embeddings (typically 384-dimensional from \texttt{sentence-transformers/all-MiniLM-L6-v2}). Captures paraphrases and synonyms that lexical matching misses.

\item[Runbook Reasoning]
The process by which an LLM agent consults a synthetic runbook (a structured troubleshooting guide for a specific fault type) to determine whether an alert is actionable, requires escalation, or can be safely aggregated as noise.

\item[Tool-Calling Agent]
An LLM agent augmented with a predefined set of functions (tools) it can invoke during reasoning: \texttt{get\_related\_alerts(service, time\_window)}, \texttt{get\_service\_dependency(service)}, and \texttt{get\_runbook(fault\_type)}. Tool outputs are fed back into the agent's context for subsequent reasoning steps.

\item[Chain-of-Thought (CoT) Prompting]
A prompting technique that instructs the LLM to generate intermediate reasoning steps before producing a final answer, improving accuracy on multi-step causal inference tasks~\cite{wei2022cot}.

\item[Fast Filter (Deduplication + Time-Window Pre-Grouping)]
A lightweight deterministic preprocessing stage that removes exact duplicate alerts and performs initial grouping by service and time window, reducing the alert volume passed to the LLM agent by an estimated 70--80\%.

\item[Synthetic Runbook]
A researcher-authored troubleshooting document for each fault type in the benchmark dataset (e.g., CPU stress, memory leak, network latency), containing symptoms, diagnosis steps, and remediation actions. Not a production SOP; used solely for experimental evaluation.

% =============================================================================
% EVALUATION METRICS
% =============================================================================

\item[Alert Reduction Ratio (ARR)]
The primary metric quantifying alert volume reduction:
\begin{equation}
\text{ARR} = \left(1 - \frac{N_{\text{groups}}}{N_{\text{raw\_alerts}}}\right) \times 100\%
\end{equation}
where $N_{\text{raw\_alerts}}$ is the number of warning alerts before aggregation and $N_{\text{groups}}$ is the number of aggregated groups presented to the engineer. Target: $\text{ARR} \ge 60\%$.

\item[Root Cause Preservation Rate (RCPR)]
The critical safety metric ensuring the root cause evidence is not lost during aggregation. It is defined at two granularities. At service granularity, which is the primary metric and spans every case that exposes logs ($M = 359$):
\begin{equation}
\text{RCPR}_{\text{svc}} = \frac{1}{M} \sum_{i=1}^{M} \mathbb{I}(s_{\text{root}}^{(i)} \in \text{PrimaryGroup}_i) \times 100\%
\end{equation}
At alert granularity, a secondary metric restricted to the $M'$ cases whose ground-truth \texttt{root\_cause.txt} indicator line is itself a warning-tier log ($M' = 4$):
\begin{equation}
\text{RCPR}_{\text{alert}} = \frac{1}{M'} \sum_{i=1}^{M'} \mathbb{I}(\text{RootCauseAlert}_i \in \text{PrimaryGroup}_i) \times 100\%
\end{equation}
where $\mathbb{I}$ is the indicator function. Target: $\text{RCPR} \ge 95\%$ (mandatory).

\item[Pairwise Grouping Precision ($P_p$)]
The fraction of alert pairs placed in the same group that truly share the same root cause:
\begin{equation}
P_p = \frac{|\{(a_x, a_y) : \text{same group} \land \text{same root cause}\}|}{|\{(a_x, a_y) : \text{same group}\}|}
\end{equation}

\item[Pairwise Grouping Recall ($R_p$)]
The fraction of alert pairs sharing the same root cause that are successfully grouped together:
\begin{equation}
R_p = \frac{|\{(a_x, a_y) : \text{same group} \land \text{same root cause}\}|}{|\{(a_x, a_y) : \text{same root cause}\}|}
\end{equation}

\item[Pairwise Grouping F1-Score ($F1_p$)]
The harmonic mean of pairwise precision and recall:
\begin{equation}
F1_p = 2 \times \frac{P_p \times R_p}{P_p + R_p}
\end{equation}
Target: $F1_p \ge 0.85$.

\item[Group Purity]
The average homogeneity of root causes within each generated group:
\begin{equation}
\text{Purity} = \frac{1}{N_{\text{groups}}} \sum_{g=1}^{N_{\text{groups}}} \frac{\max_{c} |\{a \in g : \text{root\_cause}(a) = c\}|}{|g|}
\end{equation}
Target: $\text{Purity} \ge 0.80$.

\item[Processing Latency]
The wall-clock time required to process 100 warning alerts through the aggregation pipeline, measured in milliseconds. For Arm 4, includes LLM inference time and token overhead.

% =============================================================================
% DATASET & BENCHMARK TERMS
% =============================================================================

\item[RCAEval]
A benchmark for root cause analysis in microservice systems comprising 735 failure cases across three systems (Online Boutique, Sock Shop, Train Ticket) with logs, metrics, traces, and root cause labels~\cite{pham2025rcaeval}. Licensed under MIT. Direct inspection shows that only \textbf{359 of the 735 cases ship raw log telemetry} (suite RE1 is metric-only) and that the log schema contains \texttt{timestamp}, \texttt{container\_name}, and \texttt{message} but \textbf{no severity column}.

\item[Severity Tier Assignment]
The three-tier procedure by which this work assigns the severity tier $\ell_i$, since RCAEval supplies none: \emph{Tier 1 (explicit)} parses a level token from the message text (Sock Shop, Train Ticket); \emph{Tier 2 (frequency-inferred)} flags a template as anomalous when its fault-window frequency exceeds its normal-window frequency by a ratio $R$ (Online Boutique); \emph{Tier 3 (ground-truth marker)} tags the \texttt{root\_cause.txt} indicator line as \texttt{ROOT\_CAUSE} for the eight annotated cases.

\item[LEMMA-RCA]
A large-scale multi-modal dataset for root cause analysis containing hundreds of microservice entities and four fault types (DDoS, storage failure, resource contention, noisy neighbor)~\cite{zheng2024lemma}. Licensed under CC-BY-NC-4.0 (non-commercial). Its log pipeline aggregates raw log lines into three-dimensional time series using golden-signal keywords that cover only \texttt{error}, \texttt{exception}, and \texttt{critical} --- notably \textbf{not} \texttt{warning} --- so it offers no independent warning tier.

\item[Fault Injection]
The controlled introduction of a failure (e.g., CPU stress, network latency, packet loss) into a microservice system at a known timestamp $T_{\text{inject}}$ to generate labeled failure cases for benchmarking.

\item[Normal Window]
The interval $[T_{\text{start}}, T_{\text{inject}})$ preceding fault injection, during which the system operates normally; alerts occurring here are candidates for \textit{Transient} noise. Note that the dataset fields \texttt{normal\_timesteps} and \texttt{faulty\_timesteps} are integer counts of sample points, not timestamp lists, and therefore define these windows only indirectly.

\item[Fault Window]
The interval $[T_{\text{inject}}, T_{\text{end}}]$ following fault injection, during which the system exhibits degradation; alerts here are labeled as \textit{Signal} (if from the root cause service) or \textit{Cascading} (if from a service reachable from the root cause service in the call graph).

\item[Ground Truth Root Cause Label]
The authoritative annotation of which service is the root cause of a failure case, provided by the RCAEval benchmark. Used to compute RCPR and pairwise metrics.

% =============================================================================
% BASELINE & INFRASTRUCTURE TERMS
% =============================================================================

\item[Prometheus Alertmanager]
The de facto standard alert routing and grouping component in the Cloud Native Computing Foundation (CNCF) ecosystem. Groups alerts by label matchers and time windows (\texttt{group\_wait}, \texttt{group\_interval}, \texttt{group\_by})~\cite{prometheus-alertmanager}.

\item[Group Wait (\texttt{group\_wait})]
The initial delay after the first alert in a group before sending a notification, allowing additional alerts to arrive and be grouped. Default: 30~seconds.

\item[Group Interval (\texttt{group\_interval})]
The minimum interval between subsequent notifications for the same group. Default: 5~minutes.

\item[Drain3]
An online log parsing algorithm that extracts log templates from raw log streams by clustering log messages based on token similarity, producing parameterized templates with variable placeholders~\cite{drain3}.

\item[Sentence-BERT (S-BERT)]
A modification of BERT that uses siamese and triplet network structures to derive semantically meaningful sentence embeddings, optimized for cosine similarity comparison~\cite{sbert}.

\item[HDBSCAN]
Hierarchical Density-Based Spatial Clustering of Applications with Noise; a density-based clustering algorithm that does not require specifying the number of clusters and can identify outliers (noise points)~\cite{hdbscan}.

% =============================================================================
% EXPERIMENTAL DESIGN TERMS
% =============================================================================

\item[Ablation Study]
An experimental procedure where individual components (e.g., \texttt{get\_service\_dependency} tool, fast filter, runbook) are removed from the full system to measure their individual contribution to performance metrics.

\item[Cross-System Generalization]
The ability of an aggregation method trained or tuned on one microservice system (e.g., Online Boutique) to maintain performance on unseen systems (e.g., Train Ticket) without retraining.

\item[Statistical Significance Test]
Paired Wilcoxon signed-rank test (non-parametric) used to compare metric distributions across the four arms on the same failure cases, with significance level $\alpha = 0.05$.

\item[Metric Gaming]
The phenomenon where a method optimizes a single metric (e.g., ARR) at the expense of operational utility (e.g., grouping all alerts into one group achieves ARR $\approx$ 100\% but RCPR = 0\%). Mitigated by the multi-metric framework (ARR, RCPR, F1, Purity).

% =============================================================================
% INDUSTRY STANDARDS & REFERENCES
% =============================================================================

\item[Google SRE Principles]
Foundational Site Reliability Engineering practices documented in \textit{Site Reliability Engineering} (Ch.~6, 10, 11) and \textit{The Site Reliability Workbook} (Ch.~5), emphasizing actionable alerts, maximum 2 incidents per 12-hour shift, and alerting on SLOs rather than raw metrics~\cite{google-sre-book, google-sre-workbook}.

\item[Observability Pillars]
The three canonical telemetry types: \textit{Logs} (immutable event records), \textit{Metrics} (aggregated numerical measurements), and \textit{Traces} (distributed request flow records)~\cite{grafana2025}.

\end{description}

% =============================================================================
% LOCAL BIBLIOGRAPHY MACROS (if not using main document's .bib)
% Include these in your main .bib or define via IEEEabrv.bib/IEEEfull.bib
% =============================================================================
% @string{srebook = "Site Reliability Engineering: How Google Runs Production Systems"}
% @string{sreworkbook = "The Site Reliability Workbook: Practical Ways to Implement SRE"}
% @article{pagerduty2020, ...}
% @article{grafana2025, ...}
% @inproceedings{kuang2024cola, ...}
% @inproceedings{pham2025rcaeval, ...}
% @article{zheng2024lemma, ...}
% @inproceedings{yu2025alertguardian, ...}
% @manual{prometheus-alertmanager, ...}
% @article{sbert, ...}
% @article{drain3, ...}
% @article{hdbscan, ...}
% @article{wei2022cot, ...}
% -----------------------------------------------------------------------------

\begin{description}

% =============================================================================
% CORE PROBLEM TERMS
% =============================================================================

\item[Alert Fatigue]
A cognitive state experienced by on-call engineers when the volume of alerts exceeds human cognitive capacity, leading to desensitization, missed critical alerts, and delayed incident response~\cite{pagerduty2020, grafana2025}.

\item[Alert Storm]
A sudden surge of alerts triggered by a single root cause propagating through dependent microservices, resulting in hundreds to thousands of alerts within minutes~\cite{kuang2024cola}.

\item[Noisy Alert Warning]
A warning-level alert generated by the monitoring system that is technically correct but does not require meaningful remediation action from an on-call engineer, classified into four categories: \textit{Duplicate}, \textit{Transient}, \textit{Unactionable}, and \textit{Cascading} (defined below).

\item[Early Warning Phase]
The pre-failure period during which a microservice exhibits performance degradation (e.g., increased latency, resource saturation) but has not yet violated Service Level Objectives (SLOs) or caused user-facing outages.

\item[Actionable Alert]
An alert that satisfies the Google SRE principle: ``every alert sent to an on-call engineer must require immediate human intervention''~\cite{google-sre-ch10, google-sre-ch11}.

% =============================================================================
% FOUR CATEGORIES OF NOISY ALERT WARNINGS
% =============================================================================

\item[Duplicate Warning Alert]
Multiple alerts originating from the same service with identical or near-identical message templates within a short time window ($\Delta t \le 60$~seconds), representing repeated observations of the same underlying condition rather than distinct incidents.

\item[Transient Warning Alert]
A warning alert that occurs during normal system operation (i.e., before fault injection time $T_{\text{inject}}$ in RCAEval) and resolves spontaneously without human intervention, typically caused by temporary network fluctuations or brief resource spikes.

\item[Unactionable Warning Alert]
A technically accurate warning alert for which no Standard Operating Procedure (SOP) or runbook exists, leaving the on-call engineer without a defined remediation path; the alert consumes attention but yields no operational value.

\item[Cascading Warning Alert]
A warning alert emitted by a downstream service that is a symptom of an upstream root cause; remediating the alerting service does not resolve the underlying issue, which originates in a different service (the root cause service).

% =============================================================================
% METHODOLOGY: FOUR EXPERIMENTAL ARMS
% =============================================================================

\item[Arm 1: Rule-Based Aggregation (Baseline)]
An alert grouping approach that clusters alerts based on static label matching (e.g., service name, alert name, severity) and fixed time windows (\texttt{group\_wait}, \texttt{group\_interval}), as implemented in Prometheus Alertmanager~\cite{prometheus-alertmanager}. Lacks awareness of service dependencies and semantic content.

\item[Arm 2: Semantic Similarity-Based Aggregation]
A method that parses raw logs into templates (via Drain3~\cite{drain3}), encodes templates into dense vector embeddings using Sentence-BERT~\cite{sbert}, and clusters alerts using density-based algorithms (HDBSCAN/DBSCAN) based on cosine similarity. Operates without topology knowledge.

\item[Arm 3: Temporal-Spatial Correlation Aggregation]
A topology-aware method that constructs a static service dependency graph (caller-callee adjacency matrix) from microservice architecture, then correlates alerts occurring within a sliding time window ($\Delta t \in \{30, 60, 120\}$~seconds) if a dependency path of length $k \le 2$ exists between the alerting services.

\item[Arm 4: LLM-Based Agentic Aggregation (Proposed)]
A hybrid two-stage approach: (1) a fast filter removes obvious duplicates and performs time-window pre-grouping; (2) an LLM agent (GPT-4o-mini / DeepSeek-V3) reasons over the remaining clusters using Chain-of-Thought prompting and three tools: \texttt{get\_related\_alerts}, \texttt{get\_service\_dependency}, and optionally \texttt{get\_runbook}. Limited to a maximum of two tool calls per grouping decision.

% =============================================================================
% TECHNICAL CONCEPTS & ARCHITECTURE
% =============================================================================

\item[Service Topology (Service Dependency Graph)]
A directed graph $G = (V, E)$ where vertices $V$ represent microservices and edges $E$ represent synchronous (HTTP/gRPC) or asynchronous (message queue) call relationships. Used to distinguish root cause services from downstream symptom services.

\item[Semantic Similarity]
A measure of meaning-based resemblance between two log messages, computed as the cosine similarity between their dense vector embeddings (typically 384-dimensional from \texttt{sentence-transformers/all-MiniLM-L6-v2}). Captures paraphrases and synonyms that lexical matching misses.

\item[Runbook Reasoning]
The process by which an LLM agent consults a synthetic runbook (a structured troubleshooting guide for a specific fault type) to determine whether an alert is actionable, requires escalation, or can be safely aggregated as noise.

\item[Tool-Calling Agent]
An LLM agent augmented with a predefined set of functions (tools) it can invoke during reasoning: \texttt{get\_related\_alerts(service, time\_window)}, \texttt{get\_service\_dependency(service)}, and \texttt{get\_runbook(fault\_type)}. Tool outputs are fed back into the agent's context for subsequent reasoning steps.

\item[Chain-of-Thought (CoT) Prompting]
A prompting technique that instructs the LLM to generate intermediate reasoning steps before producing a final answer, improving accuracy on multi-step causal inference tasks~\cite{wei2022cot}.

\item[Fast Filter (Deduplication + Time-Window Pre-Grouping)]
A lightweight deterministic preprocessing stage that removes exact duplicate alerts and performs initial grouping by service and time window, reducing the alert volume passed to the LLM agent by an estimated 70--80\%.

\item[Synthetic Runbook]
A researcher-authored troubleshooting document for each fault type in the benchmark dataset (e.g., CPU stress, memory leak, network latency), containing symptoms, diagnosis steps, and remediation actions. Not a production SOP; used solely for experimental evaluation.

% =============================================================================
% EVALUATION METRICS
% =============================================================================

\item[Alert Reduction Ratio (ARR)]
The primary metric quantifying alert volume reduction:
\begin{equation}
\text{ARR} = \left(1 - \frac{N_{\text{groups}}}{N_{\text{raw\_alerts}}}\right) \times 100\%
\end{equation}
where $N_{\text{raw\_alerts}}$ is the number of warning alerts before aggregation and $N_{\text{groups}}$ is the number of aggregated groups presented to the engineer. Target: $\text{ARR} \ge 60\%$.

\item[Root Cause Preservation Rate (RCPR)]
The critical safety metric ensuring the root cause evidence is not lost during aggregation. It is defined at two granularities. At service granularity, which is the primary metric and spans every case that exposes logs ($M = 359$):
\begin{equation}
\text{RCPR}_{\text{svc}} = \frac{1}{M} \sum_{i=1}^{M} \mathbb{I}(s_{\text{root}}^{(i)} \in \text{PrimaryGroup}_i) \times 100\%
\end{equation}
At alert granularity, a secondary metric restricted to the $M'$ cases whose ground-truth \texttt{root\_cause.txt} indicator line is itself a warning-tier log ($M' = 4$):
\begin{equation}
\text{RCPR}_{\text{alert}} = \frac{1}{M'} \sum_{i=1}^{M'} \mathbb{I}(\text{RootCauseAlert}_i \in \text{PrimaryGroup}_i) \times 100\%
\end{equation}
where $\mathbb{I}$ is the indicator function. Target: $\text{RCPR} \ge 95\%$ (mandatory).

\item[Pairwise Grouping Precision ($P_p$)]
The fraction of alert pairs placed in the same group that truly share the same root cause:
\begin{equation}
P_p = \frac{|\{(a_x, a_y) : \text{same group} \land \text{same root cause}\}|}{|\{(a_x, a_y) : \text{same group}\}|}
\end{equation}

\item[Pairwise Grouping Recall ($R_p$)]
The fraction of alert pairs sharing the same root cause that are successfully grouped together:
\begin{equation}
R_p = \frac{|\{(a_x, a_y) : \text{same group} \land \text{same root cause}\}|}{|\{(a_x, a_y) : \text{same root cause}\}|}
\end{equation}

\item[Pairwise Grouping F1-Score ($F1_p$)]
The harmonic mean of pairwise precision and recall:
\begin{equation}
F1_p = 2 \times \frac{P_p \times R_p}{P_p + R_p}
\end{equation}
Target: $F1_p \ge 0.85$.

\item[Group Purity]
The average homogeneity of root causes within each generated group:
\begin{equation}
\text{Purity} = \frac{1}{N_{\text{groups}}} \sum_{g=1}^{N_{\text{groups}}} \frac{\max_{c} |\{a \in g : \text{root\_cause}(a) = c\}|}{|g|}
\end{equation}
Target: $\text{Purity} \ge 0.80$.

\item[Processing Latency]
The wall-clock time required to process 100 warning alerts through the aggregation pipeline, measured in milliseconds. For Arm 4, includes LLM inference time and token overhead.

% =============================================================================
% DATASET & BENCHMARK TERMS
% =============================================================================

\item[RCAEval]
A benchmark for root cause analysis in microservice systems comprising 735 failure cases across three systems (Online Boutique, Sock Shop, Train Ticket) with logs, metrics, traces, and root cause labels~\cite{pham2025rcaeval}. Licensed under MIT. Direct inspection shows that only \textbf{359 of the 735 cases ship raw log telemetry} (suite RE1 is metric-only) and that the log schema contains \texttt{timestamp}, \texttt{container\_name}, and \texttt{message} but \textbf{no severity column}.

\item[Severity Tier Assignment]
The three-tier procedure by which this work assigns the severity tier $\ell_i$, since RCAEval supplies none: \emph{Tier 1 (explicit)} parses a level token from the message text (Sock Shop, Train Ticket); \emph{Tier 2 (frequency-inferred)} flags a template as anomalous when its fault-window frequency exceeds its normal-window frequency by a ratio $R$ (Online Boutique); \emph{Tier 3 (ground-truth marker)} tags the \texttt{root\_cause.txt} indicator line as \texttt{ROOT\_CAUSE} for the eight annotated cases.

\item[LEMMA-RCA]
A large-scale multi-modal dataset for root cause analysis containing hundreds of microservice entities and four fault types (DDoS, storage failure, resource contention, noisy neighbor)~\cite{zheng2024lemma}. Licensed under CC-BY-NC-4.0 (non-commercial). Its log pipeline aggregates raw log lines into three-dimensional time series using golden-signal keywords that cover only \texttt{error}, \texttt{exception}, and \texttt{critical} --- notably \textbf{not} \texttt{warning} --- so it offers no independent warning tier.

\item[Fault Injection]
The controlled introduction of a failure (e.g., CPU stress, network latency, packet loss) into a microservice system at a known timestamp $T_{\text{inject}}$ to generate labeled failure cases for benchmarking.

\item[Normal Window]
The interval $[T_{\text{start}}, T_{\text{inject}})$ preceding fault injection, during which the system operates normally; alerts occurring here are candidates for \textit{Transient} noise. Note that the dataset fields \texttt{normal\_timesteps} and \texttt{faulty\_timesteps} are integer counts of sample points, not timestamp lists, and therefore define these windows only indirectly.

\item[Fault Window]
The interval $[T_{\text{inject}}, T_{\text{end}}]$ following fault injection, during which the system exhibits degradation; alerts here are labeled as \textit{Signal} (if from the root cause service) or \textit{Cascading} (if from a service reachable from the root cause service in the call graph).

\item[Ground Truth Root Cause Label]
The authoritative annotation of which service is the root cause of a failure case, provided by the RCAEval benchmark. Used to compute RCPR and pairwise metrics.

% =============================================================================
% BASELINE & INFRASTRUCTURE TERMS
% =============================================================================

\item[Prometheus Alertmanager]
The de facto standard alert routing and grouping component in the Cloud Native Computing Foundation (CNCF) ecosystem. Groups alerts by label matchers and time windows (\texttt{group\_wait}, \texttt{group\_interval}, \texttt{group\_by})~\cite{prometheus-alertmanager}.

\item[Group Wait (\texttt{group\_wait})]
The initial delay after the first alert in a group before sending a notification, allowing additional alerts to arrive and be grouped. Default: 30~seconds.

\item[Group Interval (\texttt{group\_interval})]
The minimum interval between subsequent notifications for the same group. Default: 5~minutes.

\item[Drain3]
An online log parsing algorithm that extracts log templates from raw log streams by clustering log messages based on token similarity, producing parameterized templates with variable placeholders~\cite{drain3}.

\item[Sentence-BERT (S-BERT)]
A modification of BERT that uses siamese and triplet network structures to derive semantically meaningful sentence embeddings, optimized for cosine similarity comparison~\cite{sbert}.

\item[HDBSCAN]
Hierarchical Density-Based Spatial Clustering of Applications with Noise; a density-based clustering algorithm that does not require specifying the number of clusters and can identify outliers (noise points)~\cite{hdbscan}.

% =============================================================================
% EXPERIMENTAL DESIGN TERMS
% =============================================================================

\item[Ablation Study]
An experimental procedure where individual components (e.g., \texttt{get\_service\_dependency} tool, fast filter, runbook) are removed from the full system to measure their individual contribution to performance metrics.

\item[Cross-System Generalization]
The ability of an aggregation method trained or tuned on one microservice system (e.g., Online Boutique) to maintain performance on unseen systems (e.g., Train Ticket) without retraining.

\item[Statistical Significance Test]
Paired Wilcoxon signed-rank test (non-parametric) used to compare metric distributions across the four arms on the same failure cases, with significance level $\alpha = 0.05$.

\item[Metric Gaming]
The phenomenon where a method optimizes a single metric (e.g., ARR) at the expense of operational utility (e.g., grouping all alerts into one group achieves ARR $\approx$ 100\% but RCPR = 0\%). Mitigated by the multi-metric framework (ARR, RCPR, F1, Purity).

% =============================================================================
% INDUSTRY STANDARDS & REFERENCES
% =============================================================================

\item[Google SRE Principles]
Foundational Site Reliability Engineering practices documented in \textit{Site Reliability Engineering} (Ch.~6, 10, 11) and \textit{The Site Reliability Workbook} (Ch.~5), emphasizing actionable alerts, maximum 2 incidents per 12-hour shift, and alerting on SLOs rather than raw metrics~\cite{google-sre-book, google-sre-workbook}.

\item[Observability Pillars]
The three canonical telemetry types: \textit{Logs} (immutable event records), \textit{Metrics} (aggregated numerical measurements), and \textit{Traces} (distributed request flow records)~\cite{grafana2025}.

\end{description}

% =============================================================================
% LOCAL BIBLIOGRAPHY MACROS (if not using main document's .bib)
% Include these in your main .bib or define via IEEEabrv.bib/IEEEfull.bib
% =============================================================================
% @string{srebook = "Site Reliability Engineering: How Google Runs Production Systems"}
% @string{sreworkbook = "The Site Reliability Workbook: Practical Ways to Implement SRE"}
% @article{pagerduty2020, ...}
% @article{grafana2025, ...}
% @inproceedings{kuang2024cola, ...}
% @inproceedings{pham2025rcaeval, ...}
% @article{zheng2024lemma, ...}
% @inproceedings{yu2025alertguardian, ...}
% @manual{prometheus-alertmanager, ...}
% @article{sbert, ...}
% @article{drain3, ...}
% @article{hdbscan, ...}
% @article{wei2022cot, ...}
% -----------------------------------------------------------------------------

\begin{description}

% =============================================================================
% CORE PROBLEM TERMS
% =============================================================================

\item[Alert Fatigue]
A cognitive state experienced by on-call engineers when the volume of alerts exceeds human cognitive capacity, leading to desensitization, missed critical alerts, and delayed incident response~\cite{pagerduty2020, grafana2025}.

\item[Alert Storm]
A sudden surge of alerts triggered by a single root cause propagating through dependent microservices, resulting in hundreds to thousands of alerts within minutes~\cite{kuang2024cola}.

\item[Noisy Alert Warning]
A warning-level alert generated by the monitoring system that is technically correct but does not require meaningful remediation action from an on-call engineer, classified into four categories: \textit{Duplicate}, \textit{Transient}, \textit{Unactionable}, and \textit{Cascading} (defined below).

\item[Early Warning Phase]
The pre-failure period during which a microservice exhibits performance degradation (e.g., increased latency, resource saturation) but has not yet violated Service Level Objectives (SLOs) or caused user-facing outages.

\item[Actionable Alert]
An alert that satisfies the Google SRE principle: ``every alert sent to an on-call engineer must require immediate human intervention''~\cite{google-sre-ch10, google-sre-ch11}.

% =============================================================================
% FOUR CATEGORIES OF NOISY ALERT WARNINGS
% =============================================================================

\item[Duplicate Warning Alert]
Multiple alerts originating from the same service with identical or near-identical message templates within a short time window ($\Delta t \le 60$~seconds), representing repeated observations of the same underlying condition rather than distinct incidents.

\item[Transient Warning Alert]
A warning alert that occurs during normal system operation (i.e., before fault injection time $T_{\text{inject}}$ in RCAEval) and resolves spontaneously without human intervention, typically caused by temporary network fluctuations or brief resource spikes.

\item[Unactionable Warning Alert]
A technically accurate warning alert for which no Standard Operating Procedure (SOP) or runbook exists, leaving the on-call engineer without a defined remediation path; the alert consumes attention but yields no operational value.

\item[Cascading Warning Alert]
A warning alert emitted by a downstream service that is a symptom of an upstream root cause; remediating the alerting service does not resolve the underlying issue, which originates in a different service (the root cause service).

% =============================================================================
% METHODOLOGY: FOUR EXPERIMENTAL ARMS
% =============================================================================

\item[Arm 1: Rule-Based Aggregation (Baseline)]
An alert grouping approach that clusters alerts based on static label matching (e.g., service name, alert name, severity) and fixed time windows (\texttt{group\_wait}, \texttt{group\_interval}), as implemented in Prometheus Alertmanager~\cite{prometheus-alertmanager}. Lacks awareness of service dependencies and semantic content.

\item[Arm 2: Semantic Similarity-Based Aggregation]
A method that parses raw logs into templates (via Drain3~\cite{drain3}), encodes templates into dense vector embeddings using Sentence-BERT~\cite{sbert}, and clusters alerts using density-based algorithms (HDBSCAN/DBSCAN) based on cosine similarity. Operates without topology knowledge.

\item[Arm 3: Temporal-Spatial Correlation Aggregation]
A topology-aware method that constructs a static service dependency graph (caller-callee adjacency matrix) from microservice architecture, then correlates alerts occurring within a sliding time window ($\Delta t \in \{30, 60, 120\}$~seconds) if a dependency path of length $k \le 2$ exists between the alerting services.

\item[Arm 4: LLM-Based Agentic Aggregation (Proposed)]
A hybrid two-stage approach: (1) a fast filter removes obvious duplicates and performs time-window pre-grouping; (2) an LLM agent (GPT-4o-mini / DeepSeek-V3) reasons over the remaining clusters using Chain-of-Thought prompting and three tools: \texttt{get\_related\_alerts}, \texttt{get\_service\_dependency}, and optionally \texttt{get\_runbook}. Limited to a maximum of two tool calls per grouping decision.

% =============================================================================
% TECHNICAL CONCEPTS & ARCHITECTURE
% =============================================================================

\item[Service Topology (Service Dependency Graph)]
A directed graph $G = (V, E)$ where vertices $V$ represent microservices and edges $E$ represent synchronous (HTTP/gRPC) or asynchronous (message queue) call relationships. Used to distinguish root cause services from downstream symptom services.

\item[Semantic Similarity]
A measure of meaning-based resemblance between two log messages, computed as the cosine similarity between their dense vector embeddings (typically 384-dimensional from \texttt{sentence-transformers/all-MiniLM-L6-v2}). Captures paraphrases and synonyms that lexical matching misses.

\item[Runbook Reasoning]
The process by which an LLM agent consults a synthetic runbook (a structured troubleshooting guide for a specific fault type) to determine whether an alert is actionable, requires escalation, or can be safely aggregated as noise.

\item[Tool-Calling Agent]
An LLM agent augmented with a predefined set of functions (tools) it can invoke during reasoning: \texttt{get\_related\_alerts(service, time\_window)}, \texttt{get\_service\_dependency(service)}, and \texttt{get\_runbook(fault\_type)}. Tool outputs are fed back into the agent's context for subsequent reasoning steps.

\item[Chain-of-Thought (CoT) Prompting]
A prompting technique that instructs the LLM to generate intermediate reasoning steps before producing a final answer, improving accuracy on multi-step causal inference tasks~\cite{wei2022cot}.

\item[Fast Filter (Deduplication + Time-Window Pre-Grouping)]
A lightweight deterministic preprocessing stage that removes exact duplicate alerts and performs initial grouping by service and time window, reducing the alert volume passed to the LLM agent by an estimated 70--80\%.

\item[Synthetic Runbook]
A researcher-authored troubleshooting document for each fault type in the benchmark dataset (e.g., CPU stress, memory leak, network latency), containing symptoms, diagnosis steps, and remediation actions. Not a production SOP; used solely for experimental evaluation.

% =============================================================================
% EVALUATION METRICS
% =============================================================================

\item[Alert Reduction Ratio (ARR)]
The primary metric quantifying alert volume reduction:
\begin{equation}
\text{ARR} = \left(1 - \frac{N_{\text{groups}}}{N_{\text{raw\_alerts}}}\right) \times 100\%
\end{equation}
where $N_{\text{raw\_alerts}}$ is the number of warning alerts before aggregation and $N_{\text{groups}}$ is the number of aggregated groups presented to the engineer. Target: $\text{ARR} \ge 60\%$.

\item[Root Cause Preservation Rate (RCPR)]
The critical safety metric ensuring the root cause evidence is not lost during aggregation. It is defined at two granularities. At service granularity, which is the primary metric and spans every case that exposes logs ($M = 359$):
\begin{equation}
\text{RCPR}_{\text{svc}} = \frac{1}{M} \sum_{i=1}^{M} \mathbb{I}(s_{\text{root}}^{(i)} \in \text{PrimaryGroup}_i) \times 100\%
\end{equation}
At alert granularity, a secondary metric restricted to the $M'$ cases whose ground-truth \texttt{root\_cause.txt} indicator line is itself a warning-tier log ($M' = 4$):
\begin{equation}
\text{RCPR}_{\text{alert}} = \frac{1}{M'} \sum_{i=1}^{M'} \mathbb{I}(\text{RootCauseAlert}_i \in \text{PrimaryGroup}_i) \times 100\%
\end{equation}
where $\mathbb{I}$ is the indicator function. Target: $\text{RCPR} \ge 95\%$ (mandatory).

\item[Pairwise Grouping Precision ($P_p$)]
The fraction of alert pairs placed in the same group that truly share the same root cause:
\begin{equation}
P_p = \frac{|\{(a_x, a_y) : \text{same group} \land \text{same root cause}\}|}{|\{(a_x, a_y) : \text{same group}\}|}
\end{equation}

\item[Pairwise Grouping Recall ($R_p$)]
The fraction of alert pairs sharing the same root cause that are successfully grouped together:
\begin{equation}
R_p = \frac{|\{(a_x, a_y) : \text{same group} \land \text{same root cause}\}|}{|\{(a_x, a_y) : \text{same root cause}\}|}
\end{equation}

\item[Pairwise Grouping F1-Score ($F1_p$)]
The harmonic mean of pairwise precision and recall:
\begin{equation}
F1_p = 2 \times \frac{P_p \times R_p}{P_p + R_p}
\end{equation}
Target: $F1_p \ge 0.85$.

\item[Group Purity]
The average homogeneity of root causes within each generated group:
\begin{equation}
\text{Purity} = \frac{1}{N_{\text{groups}}} \sum_{g=1}^{N_{\text{groups}}} \frac{\max_{c} |\{a \in g : \text{root\_cause}(a) = c\}|}{|g|}
\end{equation}
Target: $\text{Purity} \ge 0.80$.

\item[Processing Latency]
The wall-clock time required to process 100 warning alerts through the aggregation pipeline, measured in milliseconds. For Arm 4, includes LLM inference time and token overhead.

% =============================================================================
% DATASET & BENCHMARK TERMS
% =============================================================================

\item[RCAEval]
A benchmark for root cause analysis in microservice systems comprising 735 failure cases across three systems (Online Boutique, Sock Shop, Train Ticket) with logs, metrics, traces, and root cause labels~\cite{pham2025rcaeval}. Licensed under MIT. Direct inspection shows that only \textbf{359 of the 735 cases ship raw log telemetry} (suite RE1 is metric-only) and that the log schema contains \texttt{timestamp}, \texttt{container\_name}, and \texttt{message} but \textbf{no severity column}.

\item[Severity Tier Assignment]
The three-tier procedure by which this work assigns the severity tier $\ell_i$, since RCAEval supplies none: \emph{Tier 1 (explicit)} parses a level token from the message text (Sock Shop, Train Ticket); \emph{Tier 2 (frequency-inferred)} flags a template as anomalous when its fault-window frequency exceeds its normal-window frequency by a ratio $R$ (Online Boutique); \emph{Tier 3 (ground-truth marker)} tags the \texttt{root\_cause.txt} indicator line as \texttt{ROOT\_CAUSE} for the eight annotated cases.

\item[LEMMA-RCA]
A large-scale multi-modal dataset for root cause analysis containing hundreds of microservice entities and four fault types (DDoS, storage failure, resource contention, noisy neighbor)~\cite{zheng2024lemma}. Licensed under CC-BY-NC-4.0 (non-commercial). Its log pipeline aggregates raw log lines into three-dimensional time series using golden-signal keywords that cover only \texttt{error}, \texttt{exception}, and \texttt{critical} --- notably \textbf{not} \texttt{warning} --- so it offers no independent warning tier.

\item[Fault Injection]
The controlled introduction of a failure (e.g., CPU stress, network latency, packet loss) into a microservice system at a known timestamp $T_{\text{inject}}$ to generate labeled failure cases for benchmarking.

\item[Normal Window]
The interval $[T_{\text{start}}, T_{\text{inject}})$ preceding fault injection, during which the system operates normally; alerts occurring here are candidates for \textit{Transient} noise. Note that the dataset fields \texttt{normal\_timesteps} and \texttt{faulty\_timesteps} are integer counts of sample points, not timestamp lists, and therefore define these windows only indirectly.

\item[Fault Window]
The interval $[T_{\text{inject}}, T_{\text{end}}]$ following fault injection, during which the system exhibits degradation; alerts here are labeled as \textit{Signal} (if from the root cause service) or \textit{Cascading} (if from a service reachable from the root cause service in the call graph).

\item[Ground Truth Root Cause Label]
The authoritative annotation of which service is the root cause of a failure case, provided by the RCAEval benchmark. Used to compute RCPR and pairwise metrics.

% =============================================================================
% BASELINE & INFRASTRUCTURE TERMS
% =============================================================================

\item[Prometheus Alertmanager]
The de facto standard alert routing and grouping component in the Cloud Native Computing Foundation (CNCF) ecosystem. Groups alerts by label matchers and time windows (\texttt{group\_wait}, \texttt{group\_interval}, \texttt{group\_by})~\cite{prometheus-alertmanager}.

\item[Group Wait (\texttt{group\_wait})]
The initial delay after the first alert in a group before sending a notification, allowing additional alerts to arrive and be grouped. Default: 30~seconds.

\item[Group Interval (\texttt{group\_interval})]
The minimum interval between subsequent notifications for the same group. Default: 5~minutes.

\item[Drain3]
An online log parsing algorithm that extracts log templates from raw log streams by clustering log messages based on token similarity, producing parameterized templates with variable placeholders~\cite{drain3}.

\item[Sentence-BERT (S-BERT)]
A modification of BERT that uses siamese and triplet network structures to derive semantically meaningful sentence embeddings, optimized for cosine similarity comparison~\cite{sbert}.

\item[HDBSCAN]
Hierarchical Density-Based Spatial Clustering of Applications with Noise; a density-based clustering algorithm that does not require specifying the number of clusters and can identify outliers (noise points)~\cite{hdbscan}.

% =============================================================================
% EXPERIMENTAL DESIGN TERMS
% =============================================================================

\item[Ablation Study]
An experimental procedure where individual components (e.g., \texttt{get\_service\_dependency} tool, fast filter, runbook) are removed from the full system to measure their individual contribution to performance metrics.

\item[Cross-System Generalization]
The ability of an aggregation method trained or tuned on one microservice system (e.g., Online Boutique) to maintain performance on unseen systems (e.g., Train Ticket) without retraining.

\item[Statistical Significance Test]
Paired Wilcoxon signed-rank test (non-parametric) used to compare metric distributions across the four arms on the same failure cases, with significance level $\alpha = 0.05$.

\item[Metric Gaming]
The phenomenon where a method optimizes a single metric (e.g., ARR) at the expense of operational utility (e.g., grouping all alerts into one group achieves ARR $\approx$ 100\% but RCPR = 0\%). Mitigated by the multi-metric framework (ARR, RCPR, F1, Purity).

% =============================================================================
% INDUSTRY STANDARDS & REFERENCES
% =============================================================================

\item[Google SRE Principles]
Foundational Site Reliability Engineering practices documented in \textit{Site Reliability Engineering} (Ch.~6, 10, 11) and \textit{The Site Reliability Workbook} (Ch.~5), emphasizing actionable alerts, maximum 2 incidents per 12-hour shift, and alerting on SLOs rather than raw metrics~\cite{google-sre-book, google-sre-workbook}.

\item[Observability Pillars]
The three canonical telemetry types: \textit{Logs} (immutable event records), \textit{Metrics} (aggregated numerical measurements), and \textit{Traces} (distributed request flow records)~\cite{grafana2025}.

\end{description}

% =============================================================================
% LOCAL BIBLIOGRAPHY MACROS (if not using main document's .bib)
% Include these in your main .bib or define via IEEEabrv.bib/IEEEfull.bib
% =============================================================================
% @string{srebook = "Site Reliability Engineering: How Google Runs Production Systems"}
% @string{sreworkbook = "The Site Reliability Workbook: Practical Ways to Implement SRE"}
% @article{pagerduty2020, ...}
% @article{grafana2025, ...}
% @inproceedings{kuang2024cola, ...}
% @inproceedings{pham2025rcaeval, ...}
% @article{zheng2024lemma, ...}
% @inproceedings{yu2025alertguardian, ...}
% @manual{prometheus-alertmanager, ...}
% @article{sbert, ...}
% @article{drain3, ...}
% @article{hdbscan, ...}
% @article{wei2022cot, ...}

\bibliographystyle{IEEEtran}
\bibliography{references}

\end{document}